% 本文件由 LaTeX 排版 Agent 根据有效 translation.json 自动生成。
% author=Augustine of Hippo; work_id=1506; title=论基督的恩宠与原罪

\chapter{论基督的恩宠与原罪}

% structure: p0001 source=h1 target=omitted reason=page-title-already-rendered-by-parent

第一卷\newline{}第二卷\par

\SourceRecord{Translated by Peter Holmes and Robert Ernest Wallis, and revised by Benjamin B. Warfield.；Nicene and Post-Nicene Fathers, First Series；Vol. 5.；Edited by Philip Schaff.；Buffalo, NY: Christian Literature Publishing Co.,；1887.；<http://www.newadvent.org/fathers/1506.htm>.}

% structure: p0001 source=h1 target=section reason=page-title

\section{论基督的恩宠与原罪（卷一）}

\Emph{其中他表明白拉奇在承认恩宠时缺乏诚意，因为他将恩宠置于本性、自由意志、法律和教导之中；而且，他断言恩宠仅是他所谓的意志与行动的\Quote{可能性}，而非意志与行动本身受到天主恩宠的助佑；而且，这种助佑之恩也是天主按人的功绩赐予的；他还以为人受助佑的唯一目的是能更容易地遵守诫命。奥古斯丁检验他著作中的那些段落——他自诩曾对天主的恩宠给予明确赞许之处——并指出它们如何能被解释为指向法律和教导——即指向包括在\Quote{教导}中的天主启示和基督的榜样——或指向罪过的赦免；这些段落丝毫不能证明白拉奇真正承认基督徒的恩宠，即那通过激发炽热而光辉的爱，来辅助自然禀赋与教导以行善的恩宠；奥古斯丁以请求白拉奇认真聆听他非常喜欢引用的\Person{盎博罗削}在赞美天主恩宠时极美妙的颂辞作为结语。}\par

% structure: p0003 source=h2 target=subsection reason=relative-heading-rank; base=h2

\subsection{第一章——引言。}

关于你们身体的健康，更尤其是属灵的福祉，我们感到何等巨大的喜乐，最真诚亲爱的、蒙天主钟爱的弟兄姊妹\Person{阿尔比娜}、\Person{皮尼亚努斯}和\Person{梅拉尼娅}，我们实在无法用言语表达；因此我们将这一切留给你们自己的思想和信任，好让我们现在反而来谈论你们咨询我们的问题。确实，在天主赐予我们的能力范围内，我们不得不写下这些话语，而我们的信使又正急着动身，且我们身处诸多事务之中——我在\Place{迦太基}的事务远比在任何其他地方更为繁忙。\par

% structure: p0005 source=h2 target=subsection reason=relative-heading-rank; base=h2

\subsection{第二章——白拉奇关于我们每个行动都需要恩宠的宣认之可疑性质。}

你们的来信告诉我，你们曾请求\Person{白拉奇}以书面形式表达他对所有指控的谴责；他在你们众人面前说：\Quote{我弃绝那认为或说天主的恩宠——\newline{}藉此\BibleQuote{基督耶稣来到世界拯救罪人}\BibleRef{弟前 1:15}——不仅在我们生命中的每一小时、每一刻，而且在我们每一个行动中都不是必要的人；并认为那些试图废除这恩宠的人应受永罚。}凡听到这些话，却不知道他在其著作中清楚表达的观点——不是那些他声称被以不准确形式偷走并否认的著作，而是那些他在寄往\Place{罗马}的信中提到的著作——的人，必定会以为他的看法与真理完全一致。但凡是注意到他在其中公开宣称的内容的人，就不能不以怀疑的态度看待这些陈述。因为，尽管他将那使基督来到世界拯救罪人的天主的恩宠仅视为罪过的赦免，他仍然可以将自己的话适应于这一含义，主张这种恩宠在我们生命中的每一小时、每一刻、每一个行动中都是必要的，其意义是：当我们回忆并记住过去罪过的赦免时，我们不再犯罪，这并非借助任何外在能力的供给，而是靠我们意志的力量——它在我们的每一个行动中都提醒我们，罪过的赦免赐予了我们何等益处。再者，虽然他们惯常说基督通过自己正义的生活和教导正义之事，为我们留下了榜样，从而赐予我们避免犯罪的助佑，他们也可以在这里使自己的话适应，断言恩宠对我们每一时刻、每一个行动的必要性在于：我们应在一切行为中注视主的生活榜样。但你们的忠诚使你们能清楚地看到，这样的宣认与我们现在所讨论的真正的恩宠宣认有何不同。然而，他们的模棱两可之言是多么容易将其遮蔽和伪装起来！\par

% structure: p0007 source=h2 target=subsection reason=relative-heading-rank; base=h2

\subsection{第三章——白拉奇派眼中的恩宠。}

但这有什么可奇怪的呢？因为同一位\Person{白拉奇}，在主教训导会议的记录中毫不犹豫地谴责那些说\Quote{天主的恩宠和助佑不是为每一个行动赐予的，而是在于自由意志或法律和教导}的人——在这些方面我们曾以为他已用尽了一切遁词——他也谴责那些说天主的恩宠是按我们的功绩赐予的人；然而，事实证明，在他所出版论自由意志的著作——他在寄往\Place{罗马}的信中提到了这些著作——中，他持有的无非就是他表面上谴责的那些观点。因为他将天主避免犯罪的恩宠和助佑，要么置于本性和自由意志之中，要么置于法律和教导的恩赐之中；这样当然的结果就是：每当天主帮助一个人时，就必须被认为祂通过向他启示和教导他当行之事来帮助他远离邪恶、行善，但并没有加上祂的合作和爱的激励作为额外的助佑，使他能完成他所发现的应当履行之事。\par

% structure: p0009 source=h2 target=subsection reason=relative-heading-rank; base=h2

\subsection{第四章——白拉奇的官能体系。}

在他的体系中，他提出并区分了三种官能，据他说，天主的诫命就是藉着这些官能得以成全的——\Emph{能力}、\Emph{意愿}和\Emph{行动}：他用\Quote{能力,}指使人能够成为义人的东西；用\Quote{意愿}指人愿意成为义人的东西；用\Quote{行动,}指人实际成为义人的东西。第一种，即能力，他承认是由我们本性的创造主赐给我们的；它不在我们的能力之内，我们即使违背自己的意愿也拥有它。然而，另外两种，即意愿和行动，他断言是属于我们自己的；他将它们如此严格地归于我们，以致坚持说它们仅仅来自我们自己。简言之，按照他的观点，天主的恩宠与帮助那两种他完全归于我们自己的官能（即意愿和行动）毫无关系，而只帮助那不在我们自己能力之内、来自天主的官能，即能力；仿佛属于我们自己的官能，即意愿和行动，对于离恶行善如此有效，以致不需要神圣的帮助，而我们从天主得来的官能，即能力，却是如此软弱，以致它总是得到恩宠之助的协助。\par

% structure: p0011 source=h2 target=subsection reason=relative-heading-rank; base=h2

\subsection{第五章 [IV.]——\Person{培拉久}自述官能，并引证原文。}

\Quote{我们区分三件事，}他说，\Quote{将它们按一定顺序排列。我们首先放置\OriginalTerm{能力}{ability}，其次\OriginalTerm{意愿}{volition}，第三\OriginalTerm{现实}{actuality}。\OriginalTerm{能力}{ability}置于我们的本性中，\OriginalTerm{意愿}{volition}置于我们的意志中，\OriginalTerm{现实}{actuality}置于效果中。第一，即\OriginalTerm{能力}{ability}，本当属于天主，祂将其赐予祂的受造物；另外两个，即\OriginalTerm{意愿}{volition}和\OriginalTerm{现实}{actuality}，必须归之于人，因为它们是从意志的泉源流出的。因此，对于他的愿意和做善行，赞美属于人；或者更确切地说，既属于人，也属于天主，祂赐予他意志和工作的\OriginalTerm{能力}{ability}，并经常藉着祂恩宠的帮助甚至协助这一能力。一个人能够愿意并实现任何善行，唯独来自天主。因此，这一官能可以存在，即使另外两个不存在；但后面这两个若没有前者则不能存在。我因此可以自由地没有良好的意愿或行动；但我绝不可能没有善的\OriginalTerm{能力}{capacity}。这\OriginalTerm{能力}{capacity}与生俱来，不管我愿意与否；自然在此点上从不为自己获得自由。所有这些意思将通过一两个例子变得更加清楚。我们能够用眼睛看，不是出于我们自己；但眼睛的好坏使用则是我们自己的。同样（让我通过一个一般的例子来概括一切），我们能够行、说、想任何好事，来自那赐予我们这种\OriginalTerm{能力}{ability}并也协助此\OriginalTerm{能力}{ability}的那一位；但我们真正行善事、说善言或想善念，则出于我们自己，因为我们也能够将所有这一切转向恶。因此——这一点需要经常重复，因为你们对我们的诽谤——每当我说人能够无罪生活时，我们也藉着我们承认从祂那里领受的\OriginalTerm{能力}{capacity}而赞美天主，祂将这样的\OriginalTerm{能力}{ability}赐予我们；这里没有赞美人的机会，因为当下讨论的仅仅是天主的事；问题不在于\OriginalTerm{愿意}{willing}或\OriginalTerm{实行}{effecting}，而仅仅在于那可能是什么。}\par

% structure: p0013 source=h2 target=subsection reason=relative-heading-rank; base=h2

\subsection{第六章 [V.]——\Person{培拉久}与\Person{保禄}的不同意见。}

请看，\Person{白拉奇}的整个这一教义，就是用这些话——而且别无他话——在他的维护自由意志的论文第三卷中仔细表达的，他在其中以如此精妙的方式区分了这三件事——\Quote{能力}、\Quote{意愿}和\Quote{行动}，即\Quote{能力}、\Quote{意愿}和\Quote{现实}——，以致每当我们读到或听到他承认为了避恶行善而需要天主恩宠的帮助时，无论他所说的恩宠的帮助指什么——是法律和教导，还是其他任何东西——我们都确知他所言何意；也不会因按不同于他的意思来理解而犯任何错误。因为我们不可能不知道，按照他的信念，受到天主帮助的不是我们的\Quote{意愿}也不是我们的\Quote{行动}，而仅仅是我们的\Quote{意愿和行动的能力}，三者之中只有这一项，如他所宣称的，我们是从天主得来的。仿佛天主亲自置于我们本性中的那种官能是软弱的；而另外两项，照他的意思，是我们自己的，是如此坚强、稳固且自足，以致不需要祂的任何帮助！因此，祂并不帮助我们意愿，也不帮助我们行动，而仅仅帮助我们拥有意愿和行动的可能性。然而，宗徒却持相反意见，他说：\Quote{你们要怀着恐惧战憟，成就你们自己的救恩。}\BibleRef{斐 2:12} 为使确知他们并非仅仅在能够工作（因为他们已经在本性和教导中得到了这一点），而是在实际工作中得到天主的帮助，宗徒并没有对他们说：\Quote{因为是天主在你们内工作，使你们能够}，仿佛他们已经拥有意愿和行动作为自己的资源，不需要祂在这两方面的帮助；而是说：\BibleQuote{因为是天主在你们内工作，使你们愿意并实行，为成就祂的善意}\BibleRef{斐 2:13}；或按照其他抄本，尤其是希腊文的读法：\Quote{使你们愿意并运作}。现在，试想宗徒是否因此很早以前就通过圣神预见到将出现天主恩宠的反对者；因而没有宣告天主在我们内工作的正是这两件事，即\Quote{愿意}和\Quote{运作}，而这个人却断定它们是我们自己的，仿佛它们丝毫不受天主恩宠的帮助。\par

% structure: p0015 source=h2 target=subsection reason=relative-heading-rank; base=h2

\subsection{第七章 [VI.]——\Person{白拉奇}只主张天主的帮助是为我们的\Quote{能力}。}

然而，不要让\Person{白拉奇}以这种方式欺骗粗心蒙昧的人，甚至欺骗他自己；因为在说了\Quote{因此人因愿意和行善工而应受赞扬}之后，他仿佛是纠正自己，又加了这些话：\Quote{或者更确切地说，这种赞扬属于人和\Emph{天主}。}然而，他这样表达，并不是因为他想表示尊重正确的教义，即\Quote{是天主在我们内工作，使我们愿意并实行}；而是按照他自己的说明，他为何加后一句，理由很清楚，因为他立刻补充说：\Quote{祂为了这同一意愿和工作，将\InnerQuote{能力}赐给了他。}从他前面的话可以明显看出，他将这种能力置于我们的本性中。为了不显得他对恩宠只字未提，他加了这些话：\Quote{并且祂永远借着祂恩宠的帮助，帮助这一能力本身}——注意，是\Quote{\Emph{这一能力本身},}而不是\Quote{\Emph{意愿本身},}或\Quote{\Emph{行动本身};}因为如果他这么说，他就显然不会与宗徒的教导相悖了。但他的原话是：\Quote{这一能力本身;}意思是他置于我们本性中的三种官能中的那一种。天主\Quote{永远借着祂恩宠的帮助帮助它}。结果就是，\Quote{赞扬不属于人和天主,}因为人如此愿意以致天主也以热爱的热忱激发他的意愿，或者人如此工作以致天主也与他合作——没有祂的帮助，人算什么？但他将天主与此赞扬关联起来的方式是：若非天主在创造我们时给了我们本性，使我们能够行使意愿和行动，我们既不会愿意，也不会行动。\par

% structure: p0017 source=h2 target=subsection reason=relative-heading-rank; base=h2

\subsection{第八章——依\Person{白拉奇}派，恩宠在于内在而多重的理性光照。}

关于他承认的这种自然能力，它得到天主的恩宠的帮助，但从这段文字完全不清楚他指的是哪一种恩宠，也不知道他假设我们的自然能力在多大程度上得到恩宠的帮助。但是，正如他在其他表述得更清楚、更果断的段落中一样，我们在这里也可以看出，他心目中帮助自然能力的恩宠无非就是法律和教导。\newline{}[VII.] 因为在一段话中他说：\Quote{我们被非常无知的人认为在这件事上违反了天主的恩宠，因为我们说恩宠绝不在没有我们意志的情况下在我们内完成圣德——好像天主能在他的恩宠上施加任何命令，却不向那些他施加命令的人提供他恩宠的帮助，以便人们可以更容易地通过恩宠来完成他们因自由意志而被要求做的事。}\newline{}然后，好像要解释他指的是什么恩宠，他立刻接着说了这些话：\Quote{至于我们这边，我们并不像你以为的那样，认为这种恩宠仅仅存在于法律中，而且也存在于天主的帮助中。}\newline{}现在，谁能不希望他告诉我们他要我们理解的是什么恩宠呢？事实上，我们有最强的理由希望他告诉我们，他说恩宠不仅仅存在于法律中是什么意思。\newline{}然而，当我们处于期待的悬念中，请注意，我请求你，他还有话对我们说：\Quote{天主帮助我们，}他说，\Quote{通过他的教导和启示，当他开启我们心灵的眼睛的时候；当他为我们指出未来，使我们不被当下所吞噬的时候；当他向我们揭示魔鬼的陷阱的时候；当他用多样而不可言说的天上恩宠的礼物光照我们的时候。}\newline{}然后他以一种类似赦免的方式结束他的陈述：\Quote{这人，}他问道，\Quote{说了这一切，在你看来是否是否认恩宠的人？他难道不承认人的自由意志和天主的恩宠吗？}\newline{}但到底，他还没有超出对法律和教导的赞扬；他勤勉地灌输这作为帮助我们的恩宠，从而继续他一开始的想法，当时他说：\Quote{然而，我们承认它存在于天主的帮助中。}\newline{}确实，他认为天主的帮助必须通过多种诱惑向我们推荐；通过提出教导和启示，开启心灵的眼睛，展示未来，揭露魔鬼的诡计，以及用多样而不可言说的天上恩宠的礼物光照我们的心智——这一切当然都是为了让我们学习天主的诫命和应许。而这与把天主的恩宠放在\Quote{法律和教导}中又有什么区别呢？\par

% structure: p0019 source=h2 target=subsection reason=relative-heading-rank; base=h2

\subsection{第九章 [VIII.]——法律是一回事，恩宠是另一回事。法律的效用。}

因此，很明显，他承认那种恩宠，天主通过它向我们指出并启示我们有义务做什么；但他不承认那种恩宠，天主通过它赐予并帮助我们行动，因为对法律的认识，除非伴随恩宠的帮助，反而有助于产生违犯诫命。宗徒说：\BibleQuote{在没有法律的地方，就没有违犯；} \BibleRef{罗 4:15} 又说：\BibleQuote{若不是法律说过：\InnerQuote{不可贪恋}，我就不知道贪恋。} \BibleRef{罗 7:7} 因此，法律与恩宠如此不同，以至于法律不仅无益，而且绝对有害，除非恩宠帮助它；法律的效用可以通过这一点显示出来，即它迫使所有它证明有违犯的人投奔恩宠，以获得解救和帮助来克服他们的邪恶贪欲。因为法律只是命令，而不是帮助；它发现疾病，却不医治它；不但不医治，反而加重了未得医治的病情，使得人们更迫切更焦虑地寻求恩宠的医治，正因为如此，\BibleQuote{文字叫人死，神却叫人活。} \BibleRef{格后 3:6} \BibleQuote{如果真有过能赐与生命的法律，那么正义就确实是出于法律了。} \BibleRef{迦 3:21} 然而，法律在多大程度上给予帮助，宗徒在接下来的话中告诉我们：\BibleQuote{圣经把所有的人都禁锢在罪恶之中，为使信仰耶稣基督的恩许赐与相信的人。} \BibleRef{迦 3:22} 所以宗徒说：\BibleQuote{这样，法律成了我们的启蒙师，指向基督。} \BibleRef{迦 3:24} 这件事对于骄傲的人是有益的，叫他们更坚定、更明显地\Quote{被禁锢在罪恶之下}，免得有人傲慢地试图凭自由意志如同凭自己的资源来完成自己的成义；而是为要\BibleQuote{堵住众口，使全世界都在天主面前成为有罪的。因为没有人能因法律的行为在他面前成义：因为法律只能叫人认识罪。但现在，天主的正义在法律之外已显示出来，为法律和先知所证明。} \BibleRef{罗 3:19} 那么，如果它为法律所证明，如何又说是在法律之外显示的呢？正是因为这个缘故，说的不是\Quote{没有法律而显现的}，而是\Quote{没有法律的正义}，因为这是\Quote{天主的正义}；也就是说，这正义不是我们从法律得来的，而是从天主得来的——这正义确实不是因他命令而使我们因认识它而恐惧，而是因他赐予我们而使我们因爱它而持守它——\BibleQuote{所以，凡夸耀的，应因主而夸耀。} \BibleRef{格前 1:31}\par

% structure: p0021 source=h2 target=subsection reason=relative-heading-rank; base=h2

\subsection{第十章 [IX.]——法律服务于什么目的。}

那么，此人称法律和教训为恩宠，借此帮助我们行义，究竟有何益处？因为若要它大有助益，就必须帮助我们感到自己对恩宠的需要。事实上，无人能借法律成全法律。\BibleQuote{爱是法律的满全}\BibleRef{罗 13:10}而天主的爱不是由法律倾注在我们心中的，而是由所赐予我们的圣神。\BibleRef{罗 5:5}因此，恩宠由法律所指明，为叫恩宠能成全法律。那么，对佩拉吉乌斯而言，他以不同的措辞宣称同一件事，以避免被人理解为将恩宠置于法律和教训中——即他所声称有助于我们本性之\OriginalTerm{能力}{capacity}的那恩宠，又有何益？诚然，就我所能猜测的，他之所以害怕被人如此理解，是因为他曾谴责所有那些人，他们断言天主的恩宠和帮助并非为人的单个行为而赐予，而在于人的自由、或法律和教训中。然而，他以为借着不断使用\Quote{法律和教训}这一公式，以如此多样的措辞伪装其含义，就能逃避察觉。\par

% structure: p0023 source=h2 target=subsection reason=relative-heading-rank; base=h2

\subsection{第十一章 [X.]——佩拉吉乌斯论天主如何帮助我们：\Quote{祂应许我们未来的荣耀。}}

因为在另一段中，他详细断言善志并非天主的帮助，而是出于我们自身之后，他引用宗徒书信中的问题来质问自己；他问道：\Quote{这与宗徒的话\BibleRef{斐 2:13}\InnerQuote{是天主在你们内工作，叫你们愿意，叫你们成全}如何一致？}然后，为了消除这一明显与他自己信条完全对立的权威，他立即补充说：\Quote{祂在我们内工作，使我们愿意善、愿意圣洁，是当祂借着未来荣耀的伟大和奖赏的应许，唤醒我们脱离对世俗欲望的迷恋和只顾眼前的、如同禽兽般的爱；当祂借着启示智慧，激发我们迟钝的意志渴慕天主；当祂（这是你在另一处不敢否认的）劝勉我们从事一切善事。}还有什么比这更清楚的呢？他所说的天主在我们内工作使我们愿意善的恩宠，无非就是法律和教训？因为法律和圣经的教训中应许了未来的荣耀及其巨大奖赏。启示智慧也属于教训，而教训的进一步功能是将我们的思想导向一切善事。若教训与劝勉（或更好说劝诫）之间似乎有区别，但这在\Quote{教训}的普遍术语中也能涵盖，因为教训包含在各种论述或书信中；圣经既教训也劝勉，而在教训和劝勉的过程中同样有人力的运作。然而，我们方面却愿他有朝一日承认那恩宠——在那恩宠中，不仅应许了未来的一切荣耀，而且也被信仰和盼望；不仅启示了智慧，而且也被爱慕；不仅推荐了一切善事，而且也被敦促我们接受。因为听主在圣经中应许天国的\Quote{所有的人并不都有信德}\BibleRef{得后 3:2}；也不是所有受劝前来就主的人都受劝服，主说：\BibleQuote{凡劳苦和负重担的，你们都到我跟前来。}\BibleRef{玛 11:28}但有信德的人也正是那些被劝服前来就主的人。主自己最清楚地说明了这点，祂说：\BibleQuote{除非那派遣我的父吸引他，谁也不能到我这里来。}\BibleRef{若 6:44}几节之后，当论及不信的人时，祂说：\BibleQuote{为此，我曾对你们说过：除非由父赐给他，谁也不能到我这里来。}\BibleRef{若 6:65}这就是佩拉吉乌斯应当承认的恩宠，如果他不但愿意被称为基督徒，而且愿意是基督徒的话。\par

% structure: p0025 source=h2 target=subsection reason=relative-heading-rank; base=h2

\subsection{第十二章 [XI.]——同上继续：\Quote{祂启示智慧。}}

但对于启示智慧，我该说什么呢？因为在今生，没有人能很有希望达到赐予宗徒保禄的那些伟大启示，显然不可能认为在这些启示中向他显示的除了智慧之外还有什么。然而尽管如此，他说：\BibleQuote{为了使我不要因那启示的崇高而过于高傲，故此给了我一根刺，就是撒殚的使者来攻击我，免得我过于高傲。为了这事，我三次求主使它离开我；祂对我说：我的恩宠为你足够了，因为我的力量在软弱中得以全备。}\BibleRef{格后 12:7}无疑，如果宗徒那里已经有了那不能再增加、也不能再自高自大的圆满之爱，那么他就不再需要撒殚的使者来击打他，从而抑制因大量启示而产生的过度高傲。但这高傲除了自高自大之外还有什么？关于爱，确实有最真实的说法：\BibleQuote{爱不自夸，不自大。}\BibleRef{格前 13:4}因此，这爱在这位大宗徒身上仍是日日不断增长的，只要他的\BibleQuote{内心日新又新}\BibleRef{格后 4:6}，那么无疑，当他超越了一切自夸和自大的危险时，爱才会得以成全。但那时，他的心灵仍处于一种能被大量启示所膨胀的状态，因为在爱的坚固建筑中尚未臻于成熟；他尚未到达目标，获得那他在赛程中向前直奔的奖品。\par

% structure: p0027 source=h2 target=subsection reason=relative-heading-rank; base=h2

\subsection{第十三章 [XII.]——恩宠使我们行善。}

对于不愿忍受那烦扰过程——即这自夸的性情在达到爱德最终与最高成全之前被约束——的人来说，最恰当的说法是：\BibleQuote{我的恩宠为你足够了，因为我的能力在软弱中得以成全}\BibleRef{格后 12:9}——这软弱，并非仅如这个人所认为的肉身上的软弱，而是肉身与心智的双重软弱；因为心智与那最终完满的成全相比，也是软弱的，并且为了抑制它的骄傲，也赐给了它那撒殚的使者、肉身的刺；尽管它相较于那些还不明白圣神之事的属肉体或属血气的官能，却是非常坚强的。\BibleRef{格前 2:14}因此，既然能力在软弱中得以成全，凡不自认软弱的人，便尚未走在通往成全的路上。不过，这借由软弱使能力得以成全的恩宠，引领所有按照天主旨意被预定和蒙召的人，达到最高成全与荣耀的境界。借由这恩宠，我们不仅发现当做之事，而且也实际去做我们所发现的；不仅相信当爱之事，而且也爱我们所相信的。\par

% structure: p0029 source=h2 target=subsection reason=relative-heading-rank; base=h2

\subsection{第十四章 [XII.] — 属于天主的正义与属于法律的正义}

如果这恩宠要被称为\Quote{教导}，那么至少要这样称呼：相信天主以一种不可言喻的甜美将这些教导更深、更内在地注入，不仅是通过从外部栽种和浇灌的\Emph{他们的}媒介，也同样通过祂自己在暗中运行祂自己增长的媒介——这样，祂不仅显示真理，也赋予爱。因为天主正是这样教导那些按祂旨意蒙召的人，同时赐给他们知道当做之事，以及去做他们所知道的事。因此，宗徒对得撒洛尼人这样说：\BibleQuote{论到弟兄的爱，你们不需要我写信给你们；因为你们自己受天主教导要彼此相爱。}\BibleRef{得前 4:9}然后，为证明他们已受天主教导，他补充说：\BibleQuote{你们果然向全马其顿的众弟兄行了这事。}\BibleRef{得前 4:10}仿佛你们已受天主教导的最确凿标记，就是你们实践了所受的教导。所有按天主旨意蒙召的人都是此等品性，正如先知所记：\BibleQuote{他们都将受天主教导。}然而，那学了当做之事却不去做的人，还没有按照恩宠\Quote{受天主教导}，而只是按照法律——不是按照圣神，只是按照文字。尽管有许多人因惧怕惩罚而似乎做法律所命令的事，并非出于对正义的爱；宗徒称这类正义为\Quote{属于法律的属于他自己的正义}——是一件仿佛被命令而非被赐予的东西。确实，当它被赐予时，就不被称为我们自己的正义，而是天主的正义；因为它成为我们的，惟独因我们从天主领受了它。宗徒的话是这样：\BibleQuote{为使我得以在祂内，不是有那出于法律的属于我自己的正义，而是借着基督的信德，即出于天主、基于信德的正义。}\BibleRef{斐 3:9}那么，法律与恩宠之间的差别如此之大，以致尽管法律无疑来自天主，但\Quote{出于法律}的正义并非\Quote{出于天主}，然而因恩宠而成全的正义却是\Quote{出于天主}。一个是被称为\Quote{法律的正义}，因为它是出于对法律咒诅的恐惧而行的；另一个被称为\Quote{天主的正义}，因为它是借由恩宠的恩赐而赐予的，以致它不再是一条可畏的命令，而是一条甘美的诫命，正如圣咏中的祈祷：\BibleQuote{天主，祢是良善的，因此求祢以祢的良善教导我祢的正义；}也就是说，使我不要像奴隶那样被迫在法律的恐惧下生活，而是以爱的自由喜悦地与法律相伴生活。自由人遵守诫命时，是欣然去做的。凡以这种精神学会自己责任的人，便行他所学当做的一切。\par

% structure: p0031 source=h2 target=subsection reason=relative-heading-rank; base=h2

\subsection{第十五章 [XIV.] — 被恩宠教导的人实际来到基督面前}

如今论及此等教导，主亦说：\BibleQuote{凡听见父且学习的，都到我这里来。}\BibleRef{若 6:45}因此，对于那没有来的人，不能正确地说：\Quote{他听见并学习了应当到祂那里去，但他不愿意实行所学习的。}这种说法应用于天主藉恩宠教导的方式，是绝对不合适的。因为，如果如真理所说\Quote{凡学习的都来了}，那么当然，凡不来的就没有学习。但谁能看不见人来或不来是由其意志的决定？然而，这决定可以单独存在，如果人不来；但如果他来，则不能没有帮助；这样的帮助，使他不仅知道应当做什么，而且实际地实行他所知道的。因此，当天主教导时，不是藉法律的文字，而是藉圣神的恩宠。而且，祂如此教导，使人所学习的，不仅以认知看见，而且以选择渴望，并以行动完成。因此，藉此神圣教导的方式，意志本身和行动本身都得到帮助，而不仅仅是自然的\OriginalTerm{能力}{capacity}（愿与行）。因为如果恩宠只帮助我们的这\OriginalTerm{能力}{capacity}，主就应当说：\Quote{凡听见父且学习的，\Emph{可能}到我这里来。}然而祂没有这样说，祂的话是：\Quote{凡听见父且学习的，\Emph{都}到我这里来。}现在，Pelagius将来到的可能性置于本性，甚至——如我们先前发现他试图说的——置于恩宠（无论他所谓恩宠是什么意思），他说：\Quote{藉此这能力本身得到帮助；}然而实际的来在于意志和行动。但是，可能来的人不一定实际来，除非他也愿意并行动而来。但凡由父学习的，不仅有可能来，而且\Emph{来}；在这结果中已经包含了\OriginalTerm{能力}{capacity}之\Emph{动向}，意志之\Emph{情感}，行动之\Emph{效果}。\par

% structure: p0033 source=h2 target=subsection reason=relative-heading-rank; base=h2

\subsection{第十六章 [XV.]——我们需要天主的帮助来运用我们的能力。——以视觉为例。}

他的例子有什么用处？如果它们并没有真正实现他自己使我们更清楚理解他的意义的承诺——当然，我们不必承认它们的意思，而是可以更清楚、更公开地发现他使用它们的意图和目的。他说：\Quote{我们能用眼睛看见不来自我们，但善用或滥用视力却来自我们。}对此，诗篇中有回答，诗人对天主说：\Quote{求祢使我的眼睛转离，不注视虚妄。}虽然这是就心灵之眼说的，但由此可知，就我们肉身之眼而言，可能善用或滥用：不是字面意义上的好视力（当眼睛健康时）和坏视力（当眼睛昏花时），而是道德意义上的正用（当用于帮助无助者）和邪用（当用于满足情欲）。因为虽然受助的穷人和被贪恋的女子都由这外在眼睛看见，但怜悯或情欲却出自内在之眼。那么，为何向天主祈祷说：\Quote{求祢使我的眼睛转离，不注视虚妄}？如果那事属于我们自己能力之内，而天主并不帮助意志，为何要求那事？\par

% structure: p0035 source=h2 target=subsection reason=relative-heading-rank; base=h2

\subsection{第十七章 [XVI.]——Pelagius是否有意避免公开说一切善行都来自天主？}

他说：\Quote{我们能说话来自天主，但善用或滥用言语来自我们自己。}然而，那最卓越地使用言语的人并没有这样教导我们。因为祂说：\BibleQuote{因为不是你们说话，而是你们父的圣神在你们内说话。}\BibleRef{玛 10:20}Pelagius又补充说：\Quote{同样，为了使我能通过应用一个一般情况来包括一切——我们能够行、说、想任何善事，来自那位赋予我们这能力并也帮助它的祂。}请注意，即使在这里，他重复了他以前的意思——在这三者——能力、意志、行动——中，只有能力得到帮助。然后，为了完全陈述他想要说的，他补充道：\Quote{但我们真正行善、说善言、想善念，却来自我们自己。}他忘记了他之前为了似乎修正自己的话而说过的话；因为说了\Quote{人因愿意并实行善工而应受赞美}之后，他立刻这样修改他的陈述：\Quote{或者更确切地说，这赞美既属于人，也属于赐予他这意志和工作的能力的天主。}那么，为什么他在给出例子时没有记起这个承认，至少在这些例子之后这样说：\Quote{我们能够行、说、想任何善事，来自那位赋予我们这能力并也帮助它的祂；然而，我们真正行善、说善言、想善念，却既来自我们自己，也来自祂！}？但他没有这样说。然而，如果我没有弄错，我想我看到他为什么害怕这样说。\par

% structure: p0037 source=h2 target=subsection reason=relative-heading-rank; base=h2

\subsection{第十八章 [XVII.]——他发现Pelagius如此犹豫的原因。}

因为，当他想要指出为何此事属于我们自己的能力时，他说：\Quote{因我们能将这些行为全都转向邪恶。}因此，他不敢承认这样的行为是\Quote{ \Emph{出于我们自己和出于天主},}，唯恐有人反驳他说：\Quote{如果我们行、说、思任何善事，既出于我们自己又出于天主，因为祂赋予我们这种能力，那么由此可推出，我们行、说、思恶事，也归因于我们自己与天主，因为祂也同样赋予我们中立的能力；由此得出的结论——愿天主禁止我们承认任何这样的结论——就是：正如在善行的赞美中天主与我们相连，同样祂也必须与我们分担恶行的责备。}因为祂赋予我们的那\Quote{能力}使我们既能行善也能行恶。\par

% structure: p0039 source=h2 target=subsection reason=relative-heading-rank; base=h2

\subsection{第十九章[XVIII.]——行为的两个根：爱与贪欲；各自结出自己的果实。}

关于这个\Quote{能力}，\Person{白拉奇}在其\WorkTitle{《自由意志辩护》}第一卷中写道：\Quote{现在，}他说，\Quote{天主在我们内植入了一种趋向两方面的能力。它好比一颗丰饶多产的根本，根据人的意愿，它产生并结出不同的果实；按其栽种者的选择，它要么绽放美德的花朵，要么布满罪恶的荆棘。}他几乎不顾自己所言，在此使同一根本既产生善果又产生恶果，与福音真理和宗徒的教导相悖。因为主宣布：\BibleQuote{好树不能结坏果子，坏树也不能结好果子}\BibleRef{玛 7:18}；而且当宗徒保禄说贪财是\BibleQuote{万恶的根}\BibleRef{弟前 6:10}时，他当然向我们暗示，爱可被视为诸善之根。因此，假设两棵树，一棵好，一棵坏，代表两个人，一个好，一个坏，那么好人除了有好意志外还有什么呢？即一棵有好根的树？坏人除了有坏意志外还有什么呢？即一棵有坏根的树？从这些根和树生出的果实，就是行为、言语、思想，当善时，出于善的意志，当恶时，出于恶的意志。\par

% structure: p0041 source=h2 target=subsection reason=relative-heading-rank; base=h2

\subsection{第二十章[XIX.]——人如何成为好树或坏树。}

一个人成为好树，是在他领受天主的恩宠时。因为他不是凭自己使自己由恶变善，而是出于那常是善的祂、借着祂并在祂内。而且，为了他不仅是一棵好树，还能结出好果子，他必须被同一恩宠所助佑，没有它，他不能行任何善事。因为天主亲自在好树上合作结果子：祂既通过其仆人的服务从外部浇灌和照顾它们，又由祂自己在内部使之生长\BibleRef{格前 3:7}。然而，一个人成为坏树，是在他使自己败坏时，当他背离那不变的善时；因为这种背离是恶的意志的起源。这种堕落并不导致另一种败坏的本性，而是败坏了那已被造成善的本性。然而，当这败坏被治愈后，就没有邪恶存留；因为尽管本性无疑受了损伤，但本性本身并不是瑕疵。\par

% structure: p0043 source=h2 target=subsection reason=relative-heading-rank; base=h2

\subsection{第二十一章[XX.]——爱是一切善事的根；贪欲是一切恶事的根。}

因此，我们所说的\Quote{能力}并非（如他所认为的）善恶之事的同一根源。因为作为善事之根的爱与作为恶事之根的贪欲大不相同——确实，正如美德与恶习的不同。但毫无疑问，这\Quote{能力}能容纳两者之根：因为人不仅能拥有爱，使树成为好树；他也能拥有贪欲，使树成为坏树。然而，这人类的贪欲，是一种恶习，其作者是人或人的欺骗者，而非人的创造者。它正是那\BibleQuote{肉身的贪欲，眼目的贪欲以及生活的骄傲，不是出于父，而是出于世界}\BibleRef{若一 2:16}。有谁不知道圣经的习惯用法，即以\Quote{\Emph{世界}}这一名称来指称居住在世界中的人呢？\par

% structure: p0045 source=h2 target=subsection reason=relative-heading-rank; base=h2

\subsection{第二十二章[XXI.]——爱是善的意志。}

然而，那作为德行的爱德，并非来自我们自己，而是来自天主，正如圣经所证明的：\Quote{爱德出于天主；凡有爱德的，都是生于天主，也认识天主，因为天主是爱。}\BibleRef{若一 4:7}正是基于这一爱德，才能最好地理解这段话：\Quote{凡生于天主的，就不犯罪。}\BibleRef{若一 3:9}以及这句话：\Quote{他也不能犯罪。}因为我们生于天主所凭借的爱德\Quote{不作无礼之事}，也\Quote{不图谋恶事}。\BibleRef{格前 13:5}因此，每当人犯罪时，都不是按照爱德；而是按照贪欲犯罪；遵循这样的倾向，他就不是生于天主。因为，正如已经说过的，我们所谈论的\Quote{能力}对两种根源都是开放的。所以，当圣经说：\Quote{爱德出于天主}，或更明确地说：\Quote{天主是爱}；当宗徒若望如此强调地宣告：\Quote{请看父赐给我们何等的爱情，使我们得称为天主的子女，而且我们也真是如此！}\BibleRef{若一 3:1}时，这位作者听到\Quote{天主是爱}之后，怎么还能坚持他的意见，认为我们从天主那里只领受了三者之一，即\Quote{能力}；而\Quote{善意}和\Quote{善行}却是来自我们自己？好像这善意不同于那爱德，而圣经大声宣告爱德是从天主而来，并由圣父赐给我们，好使我们成为祂的儿女。\par

% structure: p0047 source=h2 target=subsection reason=relative-heading-rank; base=h2

\subsection{第二十三章——关于恩宠赐予的根基，白拉奇的两面手法。}

或许，我们自己的先前的功绩导致这一恩赐赐予我们；正如这位作者在他写给一位圣洁贞女的作品中已经暗示的，该作品在他寄往罗马的信中被提及。因为，在引用了雅各伯宗徒的证言，他说：\Quote{你们要服从天主，抵抗魔鬼，魔鬼就会逃避你们}\BibleRef{雅 4:7}之后，他接着说：\Quote{他向我们展示了我们应该如何抵抗魔鬼，如果我们确实服从天主，并通过遵行祂的旨意而配得祂的恩宠，并在圣神的帮助下更容易地抵抗恶神。}那么，请判断他在巴勒斯坦公会议上谴责那些说天主的恩宠是按照我们的功绩赐予我们的人时，态度有多么真诚！我们还能怀疑他仍然持有这一意见，并最公开地宣扬它吗？然而，他对主教们的那份宣认怎会是真实和实在的呢？他是否已经写下了那本书，其中他最明确地声称恩宠是按照我们的功绩赐予我们的——这正是他在东方公会议上毫无保留地谴责的立场？让他坦白地承认他曾经持有这一意见，但已不再持有；这样我们就会最坦率地为他的进步而高兴。然而，事实上，除了其他指责之外，这一我们现在讨论的指责也被提出来控告他，他回答说：\Quote{无论这些是否是塞莱斯提乌斯的意见，那是那些断言如此的人的事。就我自己而言，我从未持有这样的观点；相反，我诅咒每一个持有这些观点的人。}但是，他怎么可能\Quote{从未持有这样的观点}，当时他已经写成了这部作品？或者，如果他后来写成了这部作品，他又怎能仍然\Quote{诅咒每一个持有这些观点的人}？\par

% structure: p0049 source=h2 target=subsection reason=relative-heading-rank; base=h2

\subsection{第二十四章——白拉奇将自由意志置于一切转向天主求恩宠的基础之上。}

但或许他会这样反驳我们：在他前面的这句话中，他说我们\Quote{因承行天主的旨意而堪当获得天主的恩宠}，意思是恩宠赐予那些相信并度虔诚生活的人，使他们能勇敢地抵挡诱惑者；而他们最初领受恩宠，正是为了能承行天主的旨意。那么，为防他如此反驳，请看他关于此事的另一些话语：他说：\Quote{那人急忙投奔上主，并愿受祂引导，也就是说，使自己的意志服从于天主的意志，并且如此亲密地依附上主，以致（如宗徒所说）与祂成为\InnerQuote{一个心神}\BibleRef{格前 6:17}，这一切无不出于他的\OriginalTerm{自由意志}{freedom of will}。}请注意，他宣称如此伟大的效果仅凭我们的自由意志就能达成；而且事实上，他认为我们无需天主的帮助就能依附天主：因为他如此有力地说：\Quote{无不出于他自己的自由意志。}因此，在我们无需祂帮助而依附主之后，我们甚至因这种自身的依附而堪当得到帮助。[XXIII.] 他接着说：\Quote{凡正确使用这个的}（即正确使用他的自由意志），\Quote{就完全自我奉献给天主，并彻底克己自己的意志，以致能与宗徒同说：\InnerQuote{如今活着的不再是我，而是基督在我内活着；}\BibleRef{迦 2:20}以及\InnerQuote{他将自己的心放在天主手中，使祂随意转动。}}\BibleRef{箴 21:1}天主的恩宠之助确实伟大，祂随己意转动我们的心。但照这位作者愚蠢的意见，无论帮助多大，我们都是在没有意志自由之外任何助佑的情况下，急忙投奔上主、渴望祂的引导和指引、完全将自己的意志置于祂的意志之下、并藉紧密依附而成为与祂一心的那一刻，就配得一切帮助。而所有这些浩大善行，我们（照他所说）竟单凭我们自由意志的自由就能完成；并且因这些先行的功德，我们获得了祂的恩宠，以致祂随己意转动我们的心。然而，那\Emph{那}恩宠如何不是白白赐予的呢？若是作为债务偿还，它怎能是恩宠呢？宗徒所说\Quote{这不是出于你们自己，而是天主的恩赐；不是出于行为，免得有人自夸}\BibleRef{弗 2:8}，以及\Quote{既是出于恩宠，就不是出于行为，不然恩宠就不再是恩宠了}\BibleRef{罗 11:6}，这些话怎能是真的呢？如果先有了如此有功的行为，使我们获得恩宠的赐予，那么，我再问，这些话怎能是真的？显然，在这种情况下，就没有白白赠送的礼物，只有应得报酬的回报。那么，为了找到天主的帮助，人们是否在没有天主帮助的情况下奔向天主呢？而为了在依附祂时得到天主的帮助，我们是否在没有祂帮助的情况下依附天主呢？恩宠本身还能赐予人什么更大的恩赐，或者甚至什么类似的恩赐，如果人已经不用恩宠，单凭自己自由意志的力量就能使自己与主成为一个心神呢？\par

% structure: p0051 source=h2 target=subsection reason=relative-heading-rank; base=h2

\subsection{第二十五章——天主以祂奇妙的德能在我们的心中运作我们意志的善好倾向。}

现在我要他告诉我们，那位\Place{亚述}王——他的圣洁妻子\Person{艾斯德尔}\Quote{厌恶他的床铺}——坐在王座上，穿着他所有荣耀的衣袍，全身装饰着黄金和宝石，当他抬起因愤怒而发红的脸，在辉煌中怀着野牛般的怒视看着她时，威严可畏；王后害怕了，脸色苍白，昏了过去，伏在前面使女的头上\BibleRef{艾 5:1}——我要他告诉我们，这位国王是否已经\Quote{急忙投奔上主，并愿受祂引导，使自己的意志服从于祂，并藉紧紧地依附天主而成为与祂一个心神，单凭他自由意志的力量}？他是否完全自我奉献给天主，彻底克己自己的意志，并将自己的心放在天主的手中？我想，任何人在那种状态下对国王存有这种想法，不仅愚蠢，甚至疯狂。然而，天主转变了他，将他的愤怒转为温和。然而，谁能看不出，将愤怒完全转变为温和，比当人心不为任何情感所占据，处于两者之间的冷漠中立时，使其倾向某事，要艰难得多呢？因此，让他们读和明白，观察和承认，不是通过从外面教训的法律和教导，而是通过内在运作的隐秘、奇妙、不可言喻的德能，天主在人心中运作，不仅启示真理，也赋予意志善好的倾向。\par

% structure: p0053 source=h2 target=subsection reason=relative-heading-rank; base=h2

\subsection{第二十六章——白拉奇主义的\Quote{能力}恩宠被驳斥。圣经教导在行事、言语和思想上同样需要天主的帮助。}

因此，愿\Person{白拉奇}最终停止以反对天主恩宠的辩论来欺骗自己和他人。天主的恩宠不应只因为这三者之一——即善的意志和行为的\Quote{能力}而受宣扬；也应因善的\Quote{意志}和\Quote{行为}本身而受宣扬。的确，按他的定义，这\Quote{能力}可用于两个方向；然而，我们的罪也不可因此归咎于天主，正如按他的观点，我们的善行因同一能力而归功于祂一样。因此，天主的恩宠之助之所以被主张，不仅仅是因为它帮助我们的自然能力。他必须停止说：\Quote{我们能行、说、想任何善事，是来自那位赐予我们这能力并协助这能力者；而真正行善事、说善言、想善念，则出自我们自己。}我必须重复，他必须停止说这话。因为天主不仅赐予我们能力并协助它，还更进一步在我们内工作\Quote{使你们愿意，并使你们实行}。\BibleRef{斐 2:13} 我们并非因不愿或不行善而什么都不做，而是由于没有祂的帮助。他怎能说：\Quote{我们能行善来自天主，但实际行善却出自我们自己}，当宗徒告诉我们，他为那些他所写信的人\Quote{向天主祈祷}，\Quote{使他们不做恶事，而做善事}？他的话不是\Quote{我们祈祷你们\Emph{能够}不做恶事}，而是\Quote{你们不做恶事}。他也没有说\Quote{你们\Emph{能够}行善}，而是\Quote{你们行善}。因为经上记载：\Quote{凡受天主圣神引导的，都是天主的子女}，\BibleRef{罗 8:14} 由此可知，为了行善，他们必须受那善者引导。白拉奇怎能说：\Quote{我们能善用言语来自天主，而实际善用言语则出自我们自己}，当主宣称：\Quote{因为不是你们说话，而是你们父的圣神在你们内说话}？\BibleRef{玛 10:20} 祂不是说：\Quote{不是你们自己给了自己善言的能力}，而是说：\Quote{不是你们说话}。\BibleRef{玛 10:20} 祂也没有说：\Quote{那是你们父的圣神\Emph{给了}或\Emph{曾给了你们能力}善言}，而是说：\Quote{在你们内说话}。祂指的不是\Quote{能力}的动向，而是声明合作的效果。这个自由意志的傲慢主张者怎能说：\Quote{我们能想善念来自天主，而实际想善念则出自我们自己}？他从恩宠的谦卑宣讲者那里得到回答，这位宣讲者说：\Quote{并不是我们凭自己能够思考什么，好像出于自己，而是我们的充足来自天主。}\BibleRef{格后 3:5} 注意，他没有说：\Quote{\Emph{能够}思考什么}，而是\Quote{思考什么}。\par

% structure: p0055 source=h2 target=subsection reason=relative-heading-rank; base=h2

\subsection{第二十七章 [XXVI.]——什么是真正的恩宠，以及为何赐予。功劳不先于恩宠。}

现在，连\Person{白拉奇}也应坦白承认，这恩宠在默感圣经中清楚地被陈明；他不应厚颜无耻地隐瞒他长久以来反对它的事实，而应以有益的忏悔承认它，以便圣教会不再因他的顽固坚持而受困扰，反而因他真诚的悔改而喜乐。他应区别知识和爱德，正如它们应当被区别一样；因为\Quote{知识只会使人傲慢自大，爱德才能立人。}\BibleRef{格前 8:1} 当爱德建立时，知识就不再傲慢。既然两者都是天主的恩赐（虽然一个较小，另一个更大），他不可将我们的正义高举到超过那使我们成义者应得的赞美，以至于将这两种恩赐中较小的一个归给神圣恩宠的帮助，而将较大的一个归给人的意志。如果他同意我们从天主的恩宠领受爱德，他就不应认为任何我们自己的功劳先于我们领受这恩赐。因为当我们不爱天主时，我们可能有什么功劳呢？确实，为了领受那使我们能爱的爱德，我们是在自己尚无爱德时被爱的。宗徒若望最清楚地宣告：\Quote{不是我们爱了天主，而是祂爱了我们，} \BibleRef{若一 4:10} 又说：\Quote{我们爱，因为祂先爱了我们。}\BibleRef{若一 4:19} 说得极好且真实！因为我们不可能有爱祂的资本，除非我们从祂的先爱中领受了它。如果我们没有爱，我们能做什么善事呢？如果我们有爱，我们怎能不做善事呢？虽然天主的诫命有时似乎被那些不爱祂而只畏惧祂的人遵守；但哪里没有爱，就没有善行被归算，也没有真正所谓的善行；因为\Quote{凡不出于信德的，都是罪，} \BibleRef{罗 14:23} 而\Quote{信德以爱德行使其效力。}\BibleRef{迦 5:6} 因此，天主的恩宠，藉此\Quote{祂的爱藉着所赐予我们的圣神，倾注在我们心中，} \BibleRef{罗 5:5} 必须被那愿意真实宣认的人如此宣认，以表明他毫不怀疑地相信，没有这恩宠，任何属于虔敬和真正圣善的善事都无法成就。不像那人的方式，他充分向我们显示他对恩宠的看法，当他说\Quote{赐予恩宠是为了使天主所命令的更容易履行}；这当然意味着，即使没有恩宠，天主的诫命虽较不容易，却仍能实际履行。\par

% structure: p0057 source=h2 target=subsection reason=relative-heading-rank; base=h2

\subsection{第二十八章 [XXVII.]——\Person{白拉奇}教导，没有天主恩宠的帮助，人也能抵抗撒但。}

在他写给一位圣洁贞女的书信中，有一段我已提及的文字，其中他清楚地表明了他对此事的看法；他说，我们\Quote{堪得天主的恩宠，并藉圣神的帮助\Emph{更容易地}抵抗邪灵}。为何他要插入\Quote{更容易地}一词呢？原来的意思不是已经完整：\Quote{藉圣神的帮助抵抗邪灵}吗？但谁能看不出他这一插入造成了多大的损害？当然，他想要让人以为，他急于赞扬的本性力量是如此强大，以至于即使没有圣神的帮助，也能抵抗邪灵——或许不太容易，但仍有某种程度。\par

% structure: p0059 source=h2 target=subsection reason=relative-heading-rank; base=h2

\subsection{第二十九章 [XXVIII.]——当他谈论天主的帮助时，他仅指帮助我们完成没有它我们仍能做的事。}

再者，在他\WorkTitle{意志自由辩护}的第一卷中，他说：\Quote{虽然我们内在拥有如此强大而坚定、足以抵抗犯罪的自由意志，这是我们的造物主普遍赋予人性的；然而，因着祂不可言喻的良善，我们更受到祂每日帮助的保护。}如果自由意志如此强大而坚定地抵抗犯罪，那么这种帮助有何必要呢？但这里，如同之前一样，他想要人们理解，所宣称的帮助的目的是，使人能更轻易地藉恩宠完成他仍认为没有恩宠也能——虽然不太容易，但实际上——完成的事。\par

% structure: p0061 source=h2 target=subsection reason=relative-heading-rank; base=h2

\subsection{第三十章 [XXIX.]——\Person{佩拉纠}认为为容易实行所需要的东西，实际上正是实行的必要条件。}

同样地，在同一卷的另一处，他说：\Quote{为使人们能藉恩宠更轻易地完成自由意志所命令他们去做的事。}现在，删去\Quote{\Emph{更轻易地}}一词，留下的不仅是一个完整、而且合理的含义，如果它仅简单地意指：\Quote{人们能藉恩宠完成自由意志所命令他们去做的事。}然而，加上\Quote{更轻易地}一词，则默默地暗示了即使没有天主的恩宠也能成就善工的可能性。但这样的含义是被那说\BibleQuote{离了我，你们什么也不能做}的主所禁止的。\BibleRef{若 15:5}\par

% structure: p0063 source=h2 target=subsection reason=relative-heading-rank; base=h2

\subsection{第三十一章 [XXX.]——\Person{佩拉纠}和\Person{塞莱斯提乌斯}从未真正承认恩宠。}

愿他修正这一切，免得人性软弱在如此深奥的问题上犯错后，再以魔鬼的欺骗和固执增添其错误，要么否认他确实相信的，要么坚持他曾轻率相信的，而当他回想真理之光后，发现他本不该如此相信。至于那使我们成义的恩宠——换言之，即藉着\BibleQuote{所赐给我们的圣神，将天主的爱倾注在我们心中}\BibleRef{罗 5:5}——在我有机会阅读的\Person{佩拉纠}和\Person{塞莱斯提乌斯}的著作中，我从未发现他们以应有的方式承认它。我从未在任何段落中看到他们承认\BibleQuote{应许的儿女}，关于这些儿女，宗徒如此说：\BibleQuote{肉身所生的儿女不是天主的儿女，惟独应许的儿女才算是后裔。}\BibleRef{罗 9:8}因为天主所应许的，并非我们靠自己的选择或本性的能力来实现，而是祂自己藉恩宠成就的。\par

% structure: p0065 source=h2 target=subsection reason=relative-heading-rank; base=h2

\subsection{第三十二章——为何佩拉纠派认为祈祷是必要的。\Person{佩拉纠}寄给\Person{教宗依诺森}的信函，其中阐述了他的信仰。}

现在我暂且不提\Person{塞莱斯提乌斯}的著作，或他在那些教会诉讼中提出的文件，我们已设法将全部副本连同另一封我们认为必须添加的信函寄给你们。如果你们仔细审查所有这些文件，将会发现他并未将那帮助我们避恶行善的天主恩宠置于意志的自然选择之外，而只置于法律和教导之中。他甚至断言，他们的祈祷本身就是必要的，为的是向人指明应该渴望和爱慕什么。然而，我目前可以略过这些文件不再提及；因为\Person{佩拉纠}本人最近寄了一封信和一份信仰阐述至罗马，致已故的、具有神圣记忆的\Person{教宗依诺森}，\Person{佩拉纠}当时不知其死讯。在这封信中，他说\Quote{有人试图在一些问题上诽谤他。其中之一是，他拒绝给婴儿施行圣洗圣事，并应许某些人不靠基督的救赎就能得到天国。另一个是，他谈论人避免犯罪的能力时排除了天主的帮助，如此强烈地信赖自由意志，以致拒绝了神圣恩宠的帮助。}现在，关于他关于婴儿圣洗圣事的乖僻看法（虽然他允许应为他们施行），此看法违反基督信仰和公教真理，这里不是我们准确讨论的地方，因为我们目前必须完成我们关于恩宠帮助的论著，这是我们承担的主题。让我们看看他在这封信中针对他提出的反对意见作了什么答复。我们忽略他关于对手的怨怼性抱怨，直接处理眼前主题；发现他表达如下。\par

% structure: p0067 source=h2 target=subsection reason=relative-heading-rank; base=h2

\subsection{第三十三章 [XXXI.]——\Person{佩拉纠}在恩宠主题上未声明任何不能理解为法律和教导的内容。}

\Quote{看，}他说，\Quote{这封信将如何在你的福分前为我辩白；因为我们在其中清楚而简单地声明：我们拥有一个自由意志，它对于犯罪和不犯罪是完好的；而这个自由意志在一切善工中\Emph{总是}得到天主帮助的支援。}现在，藉着主所赐给你的领悟，你看出他这话不足以解决问题。因为我们仍可追问：他所说的帮助自由意志的支援是什么？恐怕，如他惯常的，是指法律和教导。你若问他为什么用了\Quote{\Emph{总是}}这个词语，他或许会回答：因为经上记载：他昼夜默想他的法律。然后，在插入一段关于人的状况及其犯罪和不犯罪的自然能力的陈述之后，他又加上以下的话：\Quote{我们宣称这种自由意志的能力普遍存在于所有人身上——基督徒、犹太人和外邦人。所有人因本性都有同等的自由意志，但仅在基督徒身上，它才得到恩宠的支援。}我们再次追问：\Quote{什么恩宠？}他或许又会回答：\Quote{法律和基督的教导。}\par

% structure: p0069 source=h2 target=subsection reason=relative-heading-rank; base=h2

\subsection{第三十四章——白拉奇说恩宠是按人的功德赐予的。然而，功德的开始是信德；而信德是白白的恩赐，并非我们功德的回报。}

再者，无论他所说的\Quote{恩宠}是什么，他说这恩宠甚至是按基督徒的功德赐予他们的，尽管（如我以上已经提到的）当他在巴勒斯坦时，在他非常著名的自我辩护中，他谴责了持这种意见的人。以下是他原话：\Quote{在一等人中，}他说，\Quote{他们受造之善是赤身露体、毫无防卫的；}这是指那些非基督徒。然后加上其余的话：\Quote{在那些属于基督的人中，却有基督的帮助作为防卫。}你看，正如我们已经对他同一主题所评论的，这帮助究竟是什么仍然不确定。然而，他继续论及那些非基督徒说：\Quote{那些人应受审判和定罪，因为他们虽然拥有自由意志，可以藉此获得信德并配得天主的恩宠，却滥用了所赐予的自由。但这些人是应得赏报的，他们藉着正确使用自由意志而配得了主的恩宠，并遵守他的诫命。}现在很清楚，他说恩宠是按功德赐予的，无论他所说的恩宠是什么、是哪一种，但他并没有明确宣告。因为当他谈到那些善用自由意志的人应得赏报，并因此配得主的恩宠时，他实际上声称欠了他们一笔债。那么，宗徒的话\BibleQuote{蒙天主的恩宠，白白地成义}\BibleRef{罗 3:24}又成了什么？还有他另一句话\BibleQuote{你们得救是由于恩宠}\BibleRef{弗 1:8}呢？——在那里，为了防止人们以为这是出于行为，他特意加上了\BibleQuote{\Emph{藉着信德}}\BibleRef{弗 1:8}。然而，惟恐有人以为信德本身可以归于人而无关天主的恩宠，宗徒又说：\BibleQuote{而且这也不是出于你们自己，而是天主的恩赐}\BibleRef{弗 1:8}。因此，我们毫无功德地领受了那一切（按他们的说法，我们因功德所得之物）的起始——即信德本身。但如果他们坚持否认这是白白赐予我们的，那么宗徒的话\BibleQuote{按天主所分给各人信德的大小}\BibleRef{罗 12:3}又是什么意思？如果有人争辩说信德如此赐予，以致是对功德的回报，而非白白的恩赐，那么宗徒的另一句话\BibleQuote{因为你们蒙恩，不仅为相信基督，而且也要为他受苦}\BibleRef{斐 1:29}又成了什么？按宗徒的见证，二者都是恩赐——既相信基督，又为他受苦。但这些人把信德归于自由意志，以致恩宠被视为对信德的回报，而非白白的恩赐——如此恩宠便不再是恩宠，因为不是白白的，就不算恩宠。\par

% structure: p0071 source=h2 target=subsection reason=relative-heading-rank; base=h2

\subsection{第三十五章[XXXII.]——白拉奇相信婴儿在圣洗中无原罪可赦免。}

但白拉奇愿意读者从这封信转到那本陈述他信仰的书。他向你们提到过这本书，并且在其中他对一些无人质疑其观点的问题谈了不少。然而，我们只看法论及我们与他们之间争议的主题。在结束了一场他心满意足的讨论——从天主圣三的合一到肉身的复活，这些无人质问他——之后，他接着说：\Quote{我们也持守同一圣洗，我们认为这圣洗应当以与成人同样的圣事形式施行于婴儿。}那么，你们已经亲口证实你们曾听见他至少在你们面前承认了这一点。然而，他说圣洗圣事施行于儿童是\Quote{以与成人同样的言语}，当我们所追问的是事实，而非仅仅是言语时，这又有何用？更重要的一点是（正如你们所写），他亲口回答了你们的问题，说\Quote{婴儿领受圣洗是为赦免罪过。}因为他在这里也没有说\Quote{用赦免的言语}，而是承认他们是领受赦罪本身；然而尽管如此，如果你问他他认为他们被赦免的是什么罪，他会争辩说他们根本没有罪。\par

% structure: p0073 source=h2 target=subsection reason=relative-heading-rank; base=h2

\subsection{第三十六章[XXXIII.]——塞勒斯提乌斯公开宣称婴儿没有原罪。}

若在如此明确的宣认之下隐藏着相反的意义，若非塞莱斯提乌斯将其暴露，谁会相信呢？他在罗马教会诉讼程序中引用的那本书中，明确承认\Quote{婴儿也领受圣洗以得罪之赦免}，却又否认\Quote{他们有任何原罪}。但现在让我们观察白拉奇关于此事——并非关于婴儿的圣洗，而是关于天主恩宠的助佑——在他寄往罗马的那份信仰陈述中如何思考。他说：\Quote{我们如此承认自由意志，以至于我们声称自己始终需要天主的帮助。}好吧，我们再次追问，他认为我们所需要的帮助是什么；我们又一次发现歧义，因为他可能回答说，他指的是法律和基督的教导，那自然的\Quote{能力}由此得到助佑。但我们要他们承认恩宠，如同宗徒所描述的那样，他说：\BibleQuote{因为天主赐给我们的，不是胆怯之神，而是大能、爱德和健全理智之神}\BibleRef{弟后 1:7}；尽管这绝不意味着，拥有知识恩赐——由此认清自己该做什么——的人，也拥有爱德的恩宠去付诸实践。\par

% structure: p0075 source=h2 target=subsection reason=relative-heading-rank; base=h2

\subsection{第三十七章［第三十四］——白拉奇从未承认意志和行为需要天主的助佑。}

我也读过他在寄给真福教宗依诺森的书信中所提及的那些书或著作，除了一封他写给圣君士坦提乌斯主教的短信之外；但我在其中任何地方都找不到他承认这样的恩宠——即不仅助佑那\Quote{愿意和行动的自然能力}（按他的说法，即使我们既不愿善事也不行善事时，这能力仍在我们内），而且也藉着圣神的运行助佑意志和行动本身。\par

% structure: p0077 source=h2 target=subsection reason=relative-heading-rank; base=h2

\subsection{第三十八章［第三十五］——白拉奇为基督恩宠所下的定义。}

他说：\Quote{请他们阅读约十二年前我写给圣保利诺主教的信：其主题全篇约三百行，唯独宣认天主的恩宠和助佑，以及没有天主我们根本不能行任何善事。}我也读过这封信，发现他几乎只围绕本性的能力和官能展开，而天主的恩宠几乎全在于此。他确实非常简短地提及基督的恩宠，仅仅提及名称，其唯一目的似乎是为了避免完全忽略它的丑闻。然而，完全不确定他是否将基督的恩宠理解为罪过的赦免，甚或是基督的教导，包括祂生活的榜样（他在其论著的多处这样主张）；或者他是否相信，除了本性和教导之外，它还是通过炽爱之火的鼓舞而带来良好生活的帮助。\par

% structure: p0079 source=h2 target=subsection reason=relative-heading-rank; base=h2

\subsection{第三十九章［第三十六］——奥古斯丁未知的白拉奇书信。}

他说：\Quote{也请他们阅读我写给圣君士坦提乌斯主教的信，其中我——虽简短却明确——将天主的恩宠和助佑与人的自由意志结合在一起。}如同我已说过的，这封信我未读过；但若它与他所提及且我所熟悉的其它著作并无不同，则此作品对我们目前的探讨也无助益。\par

% structure: p0081 source=h2 target=subsection reason=relative-heading-rank; base=h2

\subsection{第四十章［第三十七］——白拉奇将恩宠的助佑置于教导的单纯启示中。}

再者，他说：\Quote{也请阅读我在东方写给基督的圣童贞德默特利亚的书信，他们将会发现，我们如此赞扬人的本性，却总是加上天主恩宠的助佑。}我也读了这封信；它几乎说服了我，他确实在其中承认了我们所讨论的恩宠，尽管他在这著作的许多段落中似乎确实自相矛盾。但当那些他后来为了更广泛流传而写的其它论著也到我手中时，我发现了他在何种意义上必须有意谈论恩宠——将他的信念隐藏在模糊的概括之下，却使用\Quote{恩宠}一词以减轻责难的压力并避免冒犯。因为在这著作的开头（他说：\Quote{让我们以全部热诚投入到我们面前的任务中，也不要因我们自身卑微的能力而有所疑虑；因为我们相信，我们是因母亲的信德和女儿的功德而得到助佑}），起初在我看来他承认了那种助佑我们实施个别行动的恩宠；我并没有立即注意到一个事实，即他可能将这恩宠仅仅理解为教导的启示。\par

% structure: p0083 source=h2 target=subsection reason=relative-heading-rank; base=h2

\subsection{第四十一章——白拉奇将本性的修复理解为罪过的赦免。}

在这同一著作的另一处，他说：\Quote{现在，若连没有天主的人都显明他们被天主造成怎样的人，那么基督徒——其本性已藉基督被修复到更好的状态，并且此外还得到天主恩宠的助佑——能做什么，便可想而知了。}他将这本性向更好状态的修复理解为罪过的赦免。他在此信的另一个段落中足够清楚地表明了这一点，他说：\Quote{即使那些因长期犯罪而在某种意义上变得顽固的人，也能藉着补赎得到修复。}但在此处，他甚至也可能将天主恩宠的助佑理解为教导的启示。\par

% structure: p0085 source=h2 target=subsection reason=relative-heading-rank; base=h2

\subsection{第四十二章［第三十八］——白拉奇将恩宠置于罪过的赦免和基督的榜样中。}

同样，在他这封信的另一处，他说道：\Quote{现在，如果甚至在法律之前，正如我们已经指出的，而且远在我们的主和救主来临之前，据说有些人曾经过着公义和圣洁的生活；那么，既然祂来临的光照已经使我们得以恢复，借着基督的恩宠重生为更好的人，我们能够做到这一点，岂不是更可信吗？我们这些蒙祂的血所和好并洗净，并因祂的榜样而被鼓励达到完美的义的人，岂不应该比那些生活在法律之前的人更好吗？}请注意，即使在这里，尽管措辞不同，他还是将恩宠的帮助局限于赦罪和基督的榜样。然后他补充以下的话来完成这段：\Quote{比那些生活在法律之下的人更好；按照宗徒的话，他说：\BibleQuote{罪必不能管辖你们，因为你们不在法律之下，而是在恩宠之下。}\BibleRef{罗 6:14}现在，既然我们，}他说，\Quote{已经就此谈得够多了，我想，让我们描述一位完美的贞女，她将以她行为的圣洁同时见证本性和恩宠的善，并时常以两者的德性得以温暖。}现在你应当注意，在这段话中，他也想这样结束他所说的，好让我们理解本性的善是我们受造时领受的，而恩宠的善是我们在注视和跟随基督的榜样时领受的——仿佛那些在法律之下或曾经在法律之下的人不被允许犯罪，是因为他们要么没有基督的榜样，要么不信祂。\par

% structure: p0087 source=h2 target=subsection reason=relative-heading-rank; base=h2

\subsection{第43章——白拉奇认为罪的赦免和基督的榜样足以拯救最顽固的罪人。}

确实，这就是他的意思，他的其他话语也向我们表明——这些话不包含在这部作品中，而是在他的\WorkTitle{为自由意志辩护}第三卷中，他在其中与一位对手进行讨论，对手坚持宗徒的话说：\BibleQuote{因为我所愿意的，我反而不做；}\BibleRef{罗 7:15}以及\BibleQuote{但我看见在我的肢体中另有一条法律，与我的心神的法律交战。}\BibleRef{罗 7:23}对此，他用以下的话回答：\Quote{现在，你希望我们理解这是宗徒本人说的话，而所有教会作者都断言他是以罪人、仍处于法律之下的人的身份说话——这样的人因罪恶的长期习惯，可以说被某种犯罪的必然性所束缚，并且尽管他意愿上渴望善，实际上却被拖入恶中。然而，在以一个人的身份中，}他继续说，\Quote{宗徒指出了仍在古法律之下犯罪的百姓。他宣称这个民族要借着基督从这种恶习中被拯救出来，基督首先在圣洗圣事中赦免所有信祂之人的一切罪，然后通过效法祂自己敦促他们达到完美的圣洁，并以祂自己德性的榜样战胜他们罪恶的恶习。}请注意，他认为那些在法律之下犯罪的人是如何得到帮助的：他们要通过基督的恩宠成义而被拯救，仿佛仅仅法律对他们是不够的，由于他们长期犯罪的习惯，而需要基督的某种加强；这不是借着祂的圣神注入爱德，而是借着福音中德性的教诲来默想和复制祂的榜样。现在，无论如何，这里有一个极大的需求让他明确说出他所指的恩宠是什么，因为宗徒恰恰用这些有力的话结束了构成讨论基础的段落：\BibleQuote{我真不幸啊！谁能救我脱离这该死的肉身呢？天主的恩宠，借着我们的主耶稣基督。}\BibleRef{罗 7:25}现在，当他将这恩宠不是置于祂能力的帮助中，而是置于祂作为模仿的榜样中时，我们对他还能抱有什么希望呢？因为到处都提到\Quote{恩宠}这个词，却带有模糊的普遍性？\par

% structure: p0089 source=h2 target=subsection reason=relative-heading-rank; base=h2

\subsection{第44章——白拉奇再次防备恩宠的必要性。}

然后，在他写给那位圣洁贞女的作品中（我们已经提过），有这样一段话：\Quote{让我们顺服天主，通过遵行祂的旨意，我们堪得神圣的恩宠；并且让我们借着圣神的帮助\Emph{更容易地}抵抗恶灵。}现在，从这些话中可以清楚地看出，他认为我们受到圣神恩宠的帮助，不是因为我们没有祂凭本性的纯粹能力就无法抵抗诱惑者，而是为了让我们更容易地抵抗。然而，关于这种帮助的数量和质量（无论可能是什么），我们完全可以相信他将其归结为圣神通过教导向我们启示的额外知识，这些知识我们凭本性要么不能、要么很难拥有。这些是我在他写给基督贞女的那本书中发现的细节，他在其中似乎承认了恩宠。这些细节的性质如何，你当然明白。\par

% structure: p0091 source=h2 target=subsection reason=relative-heading-rank; base=h2

\subsection{第45章——白拉奇认为应当为何目的献上祈祷。}

\Quote{让他们也读一读，}他说，\Quote{我最近出版的那本小册子，那是为了辩护自由意志而不得不早前发表的，并让他们承认，他们贬低我们否认恩宠是多么不公正，因为我们在整部著作中几乎都完全而真诚地承认了自由意志和恩宠。}这篇论文有四卷，我都读了，标注了需要思考的段落，并打算讨论它们：在读到他那封寄往罗马的书信之前，我已尽力考察了这些段落。但即使在这四卷书中，他所谓帮助我们避恶行善的恩宠，也描述得如此含糊其辞，以致他可以向追随者解释，认为恩宠所提供的帮助——为助成我们的自然能力——不外乎法律和教导。因此，我们的祈祷本身（正如他在著作中最明确地肯定）在他看来别无他用，只是为了通过神圣启示获得教导的解释，而不是为人灵提供帮助，使其以爱和行动完成所学的应行之事。事实是，他丝毫没有放弃他那体系中最明显的教义，即他提出的能力、意愿、行动这三者；他坚持只有第一个能力得享神圣帮助的持续援助，而认为意愿和行动不需要天主的帮助。此外，他所谓助成我们自然能力的帮助，他置于法律和教导之中。他承认这教导是由圣神启示或解释给我们的，因此他承认祈祷的必要性。但这法律和教导的帮助，他以为甚至在先知时代就已存在；而他所谓真正意义上的恩宠帮助，他认为仅仅在于基督的榜样。但你能清楚看到，这榜样终究属于\Quote{教导}——即作为福音宣讲给我们的教导。因此，总的结果是，仿佛为我们指出了一条道路，我们必须凭自由意志的能力行走，不需要任何他人的帮助，就能自足而不在途中昏倒或失败。甚至对于发现这道路本身，他主张单靠本性就能胜任；只是若有恩宠相助，发现会\Emph{更}容易些。\par

% structure: p0093 source=h2 target=subsection reason=relative-heading-rank; base=h2

\subsection{第四十六章〔第四十二篇〕——白拉奇声称尊重天主教作者。}

这些就是我尽力从白拉奇著作中获得的细节，只要他提及恩宠。然而，你看出，怀有上述观点的人\BibleQuote{不认识天主的义，而想建立自己的义}，\BibleRef{罗 10:3}远离\BibleQuote{从天主而来的义}，\BibleRef{斐 3:9}而非出自我们自己；他们本应在至圣的正典圣经中发现并承认这一点。然而，由于他们按自己的意思阅读这些圣经，他们自然连其中明显的事实也看不到。但愿他们能认真留意，从天主教作者的著作中至少可以学到关于天主恩宠之助的什么；因为他们自由地承认，圣经被这些作者正确理解，他们不会因过分偏爱自己的观点而忽视这些作者。请注意这位白拉基本人，在他最近发表的那篇论文中——他在自我辩护中正式提到它，即他的\WorkTitle{自由意志辩护}第三卷——赞扬了圣盎博罗削。\par

% structure: p0095 source=h2 target=subsection reason=relative-heading-rank; base=h2

\subsection{第四十七章〔第四十三篇〕——白拉奇极其赞赏盎博罗削。}

\Quote{真福盎博罗削主教，}他说，\Quote{在他的著作中，罗马信仰尤其闪耀光辉，拉丁人始终视他为作者中的花朵和荣耀，从未有敌人敢指责他的信仰或他对圣经理解的纯正。}请看他给予的赞扬的种类和分量；然而，无论他是多么神圣和有学问，他也不能与正典圣经的权威相比。如此高度赞扬盎博罗削的原因在于，白拉奇认为从后者著作中引用一段话，证明人能无罪生活。然而，这不是我们当前的问题。我们目前讨论的是帮助人避免犯罪并过圣洁生活的恩宠之助。\par

% structure: p0097 source=h2 target=subsection reason=relative-heading-rank; base=h2

\subsection{第四十八章〔第四十四篇〕——盎博罗削不与白拉奇一致。}

我确实希望他倾听这位可敬的主教，他在\WorkTitle{路加福音释义}第二卷中明确教导我们，主也合作于我们的意志。\Quote{因此你看，}他说，\Quote{因为主的能力处处与人力的努力合作，没有人能离开主建造，没有人能离开主警醒，没有人能离开主从事任何事。因此宗徒这样吩咐：\BibleQuote{你们或吃或喝，无论做什么，都要为光荣天主而行。}} \BibleRef{格前 10:31}你注意到圣盎博罗削如何甚至从人手中夺走他们惯用的说法——比如\Quote{我们从事，但天主完成}——当他在这里说\Quote{没有人能离开主从事任何事}时。同样，在同书的第六卷中，论及某人有两个债户时，他说：\Quote{按人的看法，欠得最多的人也许罪更大。但主的仁慈改变了情况，以致欠得最多的人爱得也最多，如果他获得了恩宠的话。}请看这位公教博士最清楚地宣告，促使每个人更深爱的爱本身属于恩宠的恩赐。\par

% structure: p0099 source=h2 target=subsection reason=relative-heading-rank; base=h2

\subsection{第四十九章〔第四十五篇〕——盎博罗削教导，基督以何等目光转身注视伯多禄。}

诚然，悔改本身无疑是一种意志的行为，但它是藉主的仁慈和帮助而付诸行动的；真福安波罗修在同一著作的第九卷中以下文证实了这一点。他说：\Quote{好的泪水能洗去罪恶。主最终转面垂视的人才会哭泣。伯多禄先否认了主，却没有哭，因为主还没有转面垂视他。他第二次否认主，仍没有哭，因为主还没有转面垂视他。他第三次也否认了主，耶稣转面垂视，他便痛哭流泪了。}愿这些人阅读福音；愿他们思考：那时主耶稣在内室中，正在大司祭前受审；而宗徒伯多禄在外面，\BibleRef{玛 26:69}又在庭院下方，\BibleRef{谷 14:66}时而与仆役们坐在火旁，\BibleRef{路 22:55}时而站着，\BibleRef{若 18:16}正如福音作者们最精确且一致的叙述所显示的。因此，不能说主是用肉身的眼睛，藉可见且明显的劝诫转面垂视了他。那么，经上所记载的：\Quote{主转面垂视伯多禄}\BibleRef{路 22:61}这一句话，是在内心发生的；它是在心灵中、在意志中造成的。主仁慈地、静默而隐秘地临近，触摸了心，唤起了对过去的记忆，用内在的恩宠探访了伯多禄，激发并使他内在的情感流露为外在的泪水。请看，天主是如何以祂的帮助临于我们的意志和行动；请看祂如何\Quote{在我们内运作，使我们愿意并实行}。\par

% structure: p0101 source=h2 target=subsection reason=relative-heading-rank; base=h2

\subsection{第50章——安波罗修教导众人需要天主的帮助。}

在同一书中，同一圣安波罗修又说：\Quote{如今，如果伯多禄跌倒了——他曾说：\InnerQuote{即便众人跌倒，我却总不跌倒。}——还有谁应当正确自信于自己呢？达味确实，因为他曾说过：\InnerQuote{我在顺遂中曾说：我永不动摇。}——他承认自己的自信对他多么有害：\InnerQuote{你转面不看我，}他说，\InnerQuote{我就惊慌不安。}}白拉奇应当听从如此杰出之人的教导，并追随他的信仰，因为他已经称赞了他的教导和信仰。让他谦卑地聆听，忠信地追随；不要再固执自大，以免丧亡。为什么白拉奇宁愿沉入那片大海，而伯多禄却被磐石拯救了出来？\par

% structure: p0103 source=h2 target=subsection reason=relative-heading-rank; base=h2

\subsection{第51章[XLVI.]——安波罗修教导：是天主为人类所做，而白拉奇归功于自由意志。}

也请这位敬虔的主教，在同一著作的第六卷中说：\Quote{他们不愿接待他的原因，福音作者自己以这些话指明：\InnerQuote{因为他的面容好像是往耶路撒冷去。}\BibleRef{路 9:53}但祂的门徒十分渴望祂能被撒玛黎雅人所接待。然而，天主召叫凡祂所愿意的人，并使凡祂所愿意的人成为虔诚的。}这位天主之人从天主恩宠的泉源中汲取了何等明智的洞察力！他说：\Quote{天主召叫凡祂所愿意的人，并使凡祂所愿意的人成为虔诚的。}请看，这不正是先知所宣称的：\Quote{我要恩待谁，就恩待谁；我要怜悯谁，就怜悯谁}\BibleRef{出 33:19}以及宗徒由此推论：\Quote{这样看来，那愿意的，那奔跑的，都不算什么，而是那显示仁慈的天主。}\BibleRef{罗 9:16}如今，当我们这个时代的新人物（指白拉奇）说：\Quote{凡天主所愿意的人，祂就召叫；凡祂所愿意的人，祂就使他成为虔诚的。}——有谁敢强辩说，那个人还不虔诚，\Quote{他急速投奔上主，渴望被祂引领，并将自己的意志依附于天主的意志；而且，他如此紧密地依附于主，以致成为（如宗徒所说）与主\InnerQuote{同一心神}的人？}然而，这整个\Quote{虔诚的人}的作为是多么伟大，但白拉奇却主张\Quote{这只是由自由意志完成的}。然而，他所极力赞扬的那位真福安波罗修反对他说：\Quote{主天主召叫凡祂所愿意的人，并使凡祂所愿意的人成为虔诚的。}因此，是天主使凡祂所愿意的人成为虔诚的，以便他能\Quote{急速投奔上主，渴望被祂引领，并将自己的意志依附于天主的意志，并如此紧密地依附于主，以致成为（如宗徒所说）与主\InnerQuote{同一心神}的人}；而这一切，除了虔诚的人之外无人能做。那么，有谁会做如此多的事，除非是天主使他做的？\par

% structure: p0105 source=h2 target=subsection reason=relative-heading-rank; base=h2

\subsection{第52章[XLVII.]——如果白拉奇同意安波罗修，奥古斯丁就不再与他争论。}

但是，由于关于自由意志和天主恩宠的讨论在区分上如此困难，以至于当自由意志被维护时，天主的恩宠似乎被否认；而当恩宠被肯定时，自由意志又被认为已被消除——\Person{白拉奇}可以陷入这种模糊的阴影中，以至于宣称他同意我们引用\Person{圣安博罗削}的一切话，并声明这也是他的一贯看法；并试图这样解释每一点，使人们以为他的意见与安博罗削相当一致。因此，就天主帮助和恩宠的问题而言，请你注意他所清楚区分的三件事，即\Quote{能力}、\Quote{意志}和\Quote{实行}，也就是\Quote{能力}、\Quote{意愿}和\Quote{行动}。因此，如果他同意我们的观点，那么不仅是在人里面的\Quote{能力}——即使他既不愿也不实行善——而是连意愿和实行，也就是我们愿意行善和实际行善——这些东西在人里面除非他有了好意志并正确行动，否则是不存在的——：如果，我再说一遍，他这样同意与我们一同主张，即使是意愿和实行也由天主所助，并且如此受助，以至于没有这种帮助，我们既不能意愿也不能实行任何善事；如果他相信这正是藉着我们的主耶稣基督而来的天主的恩宠，这恩宠使我们因他的义而非我们自己的义成为义，以致我们真正的义是从他而来的——那么，就我所能判断的，我们之间关于天主恩宠的帮助将不再有任何争议。\par

% structure: p0107 source=h2 target=subsection reason=relative-heading-rank; base=h2

\subsection{第五十三章 [第四十八章]——在何种意义上可以说有些人今生无罪}

但关于他特别引用圣安博罗削以表赞同的那一点——因为他在这位作者著作中发现，从他对\Person{匝加利亚}和\Person{依撒伯尔}的赞扬来看，一个人有可能在此生无罪；虽然如果天主愿意，这是不可否认的，因为在天主一切都是可能的，但他应当更仔细地考虑这\Emph{在何种意义}上说的。就我所能见到的，这句话是根据某种行为标准而说的，这种标准在人看来是值得赞赏和表扬的，没有人能公正地质疑它以便提出指控或责备。据说匝加利亚和他的妻子依撒伯尔在天主面前保持了这样的标准，没有别的原因，只是因为他们在这标准中行走，从未以任何伪善欺骗人；而是如同他们以其真诚显现在人面前，他们也在天主面前被认识。然而，这句话并非针对那种完全的义德状态，即我们将来要在其中真正绝对地生活在无瑕纯洁的情形中。事实上，宗徒\Person{保禄}曾告诉我们，在法律所要求的义德上，他是\Quote{无可指摘的}\BibleRef{斐 3:6}；匝加利亚也是本着同样的法律过无可指摘的生活。然而，宗徒将这义德视为\Quote{粪土}和\Quote{损失}\BibleRef{斐 3:8}，与我们所希望的义德相比，我们应当\Quote{饥渴}\BibleRef{玛 5:6}，以便将来当我们凭信德享有它时，得以满足于其荣观，因为\Quote{义人因信德而生活}\BibleRef{罗 1:17}。\par

% structure: p0109 source=h2 target=subsection reason=relative-heading-rank; base=h2

\subsection{第五十四章 [第四十九章]——安博罗削教导说，在这世界上无人无罪}

最后，请他好好注意他可敬的主教，当他讲解依撒意亚先知书时说：\Quote{在\Emph{这世界}上，没有人能无罪。}没有人能声称，用\Quote{在这世界}这个短语，他仅仅是指爱这世界。因为他谈及宗徒所说的话：\Quote{我们的家乡在天上}\BibleRef{斐 3:20}；这位杰出主教在阐明这些话语的意义时，这样表达：\Quote{现在宗徒说，许多人在现世生活中对自身而言是完美的，但如果着眼于真正的完美，他们决不能被认为是完美。因为他自己说：\InnerQuote{我们现在藉着镜子观看，模糊不清，到那时就要面对面了；我现在所认识的，只是局部的，到那时就全认识，如同我被认识一样。}\BibleRef{格前 13:13} 因此，有些人在现世是无瑕的，有些人在天国里将是无瑕的；当然，如果你仔细推敲，没有人能是无瑕的，因为没有人是无罪的。}因此，\Person{白拉奇}用来支持他自己意见的圣安博罗削的那段话，要么是有条件地写的，虽有可能但未精确表述；要么如果这位圣洁谦卑的作者确实认为匝加利亚和依撒伯尔按照最高且绝对完美的义德生活，无法增加或补充，那么他在更仔细的检查后，当然修正了自己的意见。\par

% structure: p0111 source=h2 target=subsection reason=relative-heading-rank; base=h2

\subsection{第五十五章 [第五十章]——安博罗削见证说，完全纯洁对于人性是不可能的}

此外，他应当仔细注意，在他所引用的同一段语境中，他从\Person{盎博罗削}那里引用了似乎对其目的十分满意的那段话，他也说了这句话：\Quote{从起初就无玷，为人的本性是不可能的。}在这句话中，可敬的\Person{盎博罗削}确实将软弱和虚弱归因于那本性的\Quote{能力}，\Person{白拉奇}却拒绝忠实地承认这种能力已被罪败坏，因而自夸地颂扬它。毫无疑问，这违反了这个人的意愿和倾向，尽管它并不抵触宗徒的真实宣认，宗徒说：\BibleQuote{我们从前也都在本性上为义怒之子，和别人一样。}\BibleRef{弗 2:3}因为通过第一个人的罪，这罪来自他的自由意志，我们的本性变得败坏和毁灭；唯有天主的恩宠，通过祂——祂是天主与人之间的中保，也是我们全能的医生——救助它。如今，由于我们在论述天主恩宠对我们的成义的助佑方面已经将此著作拉得过长——天主藉此恩宠使一切有助于那些爱祂的人，\BibleRef{罗 8:28}以及祂先爱了我们，\BibleRef{若一 4:19}——将恩赐给他们，好从他们那里领受：我们必须按主的允准，开始另一篇关于罪的论述，罪藉着一个人进入了世界，并带着死亡，如此便传给了众人，\BibleRef{罗 5:12}我们要提出看来必要且充分的论述，以反对那些对此处所述的真理爆发出猛烈和公开错误的人。\par

\SourceRecord{Translated by Peter Holmes and Robert Ernest Wallis, and revised by Benjamin B. Warfield.；Nicene and Post-Nicene Fathers, First Series；Vol. 5.；Edited by Philip Schaff.；Buffalo, NY: Christian Literature Publishing Co.,；1887.；<http://www.newadvent.org/fathers/15061.htm>.}

% structure: p0001 source=h1 target=section reason=page-title

\section{论基督的恩宠与原罪（第二卷）}

\Emph{奥古斯丁在此指出，佩拉纠在有关原罪与婴儿圣洗的问题上，与他的追随者塞莱斯提乌斯实际上毫无区别：塞莱斯提乌斯拒绝承认原罪，甚至敢在公众面前否认这一教义，因此在主教们的审判中被定罪——先是在迦太基，后来又在罗马。因为这个问题并不像这些异端者所认为的那样，是一个可以在信仰上无危险地犯错误的问题。实际上，他们的异端所针对的正是基督教信仰的根本。之后，他驳斥了所有那些坚持认为原罪教义贬低了婚姻祝福、并侮辱了天主本身——那位经由婚姻而使人诞生的造物主——的人。}\par

% structure: p0003 source=h2 target=subsection reason=relative-heading-rank; base=h2

\subsection{第一章——审慎对待佩拉纠关于婴儿圣洗的论述之必要性}

其次，我恳请你们仔细观察，在这种婴儿圣洗的问题上，你们应当带着何等审慎的态度去聆听这类人的言论：他们不敢公开否认这幼小年龄需要重生之洗和罪赦，惟恐基督徒的耳朵不愿听；然而他们却固执地坚持并鼓吹他们的意见，即肉身的生育并不承担人类第一次的罪债，尽管他们似乎允许婴儿为得赦罪而受圣洗。确实，你们在信中已告知我，你们曾亲耳听到佩拉纠在你们面前宣读他书中的那段话，那本书他声称也已送到罗马；他们坚持认为：\Quote{婴儿应当以与成年人相同的圣事词语格式受圣洗。}谁在听到这一陈述后，还会认为有必要就此问题提出任何质疑呢？或者，如果有人质疑，在读到他们明确的主张（即他们否认婴儿继承原罪，并坚持所有人生来毫无瑕疵）之前，谁不会认为他是过分吹毛求疵呢？\par

% structure: p0005 source=h2 target=subsection reason=relative-heading-rank; base=h2

\subsection{第二章——塞莱斯提乌斯在迦太基受审时拒绝谴责其错误；他提交给佐西默斯的书面声明}

塞莱斯提乌斯确实以较少克制地坚持这种错误教义。他如此放肆，以至于在迦太基主教面前受审时，竟拒绝谴责那些说\Quote{亚当的罪仅损害了亚当本人，而非人类全体；并且婴儿出生时的状况与亚当在堕落前的状况相同}的人。此外，在他提交给罗马最蒙福的教宗佐西默斯的书面声明中，他特别明确地宣称：\Quote{原罪不束缚任何婴儿。}关于迦太基的教会诉讼，我们摘录以下他的言论。\par

% structure: p0007 source=h2 target=subsection reason=relative-heading-rank; base=h2

\subsection{第三章——迦太基公会议起诉塞莱斯提乌斯的部分记录}

主教奥雷留说：\Quote{让以下内容被宣读。}随即，宣读的内容是：\Quote{亚当的罪仅损害了亚当本人，而非人类全体；}然后，在宣读之后，塞莱斯提乌斯说：\Quote{我说过，我对罪的传袭存有疑问，但愿意顺从任何蒙天主赐予知识恩宠的人；因为我甚至从那些被任命为天主教会司铎的人那里听到了不同的意见。}执事保利努斯说：\Quote{告诉我们他们的名字。}塞莱斯提乌斯回答：\Quote{圣司铎鲁菲努斯，他曾与圣潘玛丘一同住在罗马。我听到他宣称没有罪的传袭。}执事保利努斯随后问：\Quote{还有别人吗？}塞莱斯提乌斯回答：\Quote{我听更多人也这么说。}执事保利努斯反问道：\Quote{告诉我们他们的名字。}塞莱斯提乌斯说：\Quote{一个司铎还不够吗？}然后，在另一处我们读到：主教奥雷留说：\Quote{宣读剩下的指控。}随即宣读的是：\Quote{婴儿出生时的状况与亚当在堕落前的状况相同；}他们一直读完先前提交的简短指控的结尾。[IV.]主教奥雷留问道：\Quote{塞莱斯提乌斯，你是否曾像执事保利努斯所说的那样，教导过婴儿出生时的状况与亚当在堕落前的状况相同？}塞莱斯提乌斯回答：\Quote{请他也解释一下当他说到\InnerQuote{\Emph{堕落之前}}时的意思。}执事保利努斯随后说：\Quote{你方面是否否认曾教导过这个教义？两者必居其一：他必须要么说他从未如此教导，要么就必须现在谴责这个意见。}塞莱斯提乌斯反驳道：\Quote{我已经说过，请他解释他所说的\InnerQuote{\Emph{堕落之前}}这句话。}执事保利努斯随后说：\Quote{你必须否认曾教导过这个。}主教奥雷留说：\Quote{我问，我该从这人的顽固中得出什么结论；我的肯定陈述是，尽管亚当最初在乐园中被造时被说成是不朽的，但他后来因违命而变得可朽。你是否这样说，保利努斯兄弟？}执事保利努斯回答：\Quote{我是的，主上。}然后主教奥雷留说：\Quote{关于现今婴儿在受圣洗前的状况，执事保利努斯希望知道，是否与亚当在堕落前的状况相同；以及它是否从它所由生的同一罪恶根源中沾染了违命的罪责？}执事保利努斯问：\Quote{让他否认是否曾教导过这个。}塞莱斯提乌斯回答：\Quote{关于罪的传袭，我已经断言，我听到许多在天主教会中有地位的人完全否认它；而另一方面，其他人则肯定它：这可以公平地视为一个有待探究的问题，而非异端。我一直坚持婴儿需要圣洗，并且应当受圣洗。他还想要什么？}\par

% structure: p0009 source=h2 target=subsection reason=relative-heading-rank; base=h2

\subsection{第四章——塞莱斯提乌斯让步承认婴儿圣洗，但不肯定原罪}

当然，你们看到塞莱斯提乌斯在此对婴儿圣洗的让步只是如此，以至于他不愿承认第一个人的罪——那在重生之水中被洗去的——传递给他们，尽管同时他也不敢否认这一点；并且由于这种怀疑，他拒绝谴责那些主张\Quote{亚当的罪只伤害了他自己，而非整个人类}以及\Quote{婴儿出生时的状态与亚当在犯罪之前的状态相同}的人。\par

% structure: p0011 source=h2 target=subsection reason=relative-heading-rank; base=h2

\subsection{第五章 [V.]——塞莱斯提乌斯在罗马审讯中提交的书。}

但在他在罗马出版并在教会面前的审讯中提交的那本书中，他在这个问题上如此说话，以至于显示他实际上相信他曾自称怀疑的东西。因为他的话如下：\Quote{然而，婴儿应当为罪的赦免而受圣洗，按照普世教会的规则和福音的含义，我们承认。因为主已经确定，天国只应赐给受圣洗的人；\BibleRef{若 3:5}而既然自然的力量不拥有它，它必须由恩宠的恩赐所赐予。}现在，如果他在这个问题上别处没有说什么，谁不会认为他承认甚至在婴儿的圣洗中原罪的赦免，通过说他们应当为罪的赦免而受圣洗？因此，你在信中所说的这一点，即佩拉纠对你的回答是这样的：\Quote{婴儿与成人使用相同的圣事格式的词语受圣洗，}并且你很高兴听到你渴望听到的同样事情，然而你更倾向于就他的话与我们商议。\par

% structure: p0013 source=h2 target=subsection reason=relative-heading-rank; base=h2

\subsection{第六章 [VI.]——门徒塞莱斯提乌斯在此作品中比他的师傅更大胆。}

那么，请仔细观察塞莱斯提乌斯如此公开地提出的观点，你就会发现佩拉纠在你面前隐藏了多少。塞莱斯提乌斯接着如下说道：\Quote{然而，我们必须为罪的赦免而为婴儿付圣洗，这并不是因为我们似乎要肯定罪是由遗传而来的。这与公教的意义非常相悖，因为罪不是与人生俱来的——而是由人随后犯下的，因为罪被证明是意志的过错，而非本性的过错。因此，承认这一点是恰当的，以免我们似乎要制造不同的圣洗；此外，有必要提出这一初步的防范，免得因这奥迹的机缘，恶——在对造物主的贬损下——被说成是由本性传给人的，而在此之前它还未被人所犯。}现在，佩拉纠或者害怕或者羞于在你面前承认这是他自己的意见；尽管他的门徒在宗座面前公开承认这是他的意见时，既没有感到不安也没有脸红，没有任何模糊的遁词。\par

% structure: p0015 source=h2 target=subsection reason=relative-heading-rank; base=h2

\subsection{第七章——教宗佐西姆仁慈地宽恕了他。}

然而，主持此座的主教，看到他像疯子一样在如此大的狂妄中狂奔，出于极大的同情，为了那人的悔改——如果可能的话——宁愿通过从他那里引出他对自订问题的回答来紧紧束缚他，而不是通过严厉谴责的打击将他推下他看来已准备好要跌落的悬崖。我刻意说\Quote{他看来已准备好要跌落}，而不是\Quote{他已经跌落了}，因为他已经在这同一本书中，有意地为这类问题作了预备，说了以下的话：\Quote{如果碰巧任何出于无知的错误已经悄悄潜入我们人类，愿它被你的决定性判决纠正。}\par

% structure: p0017 source=h2 target=subsection reason=relative-heading-rank; base=h2

\subsection{第八章 [VII.]——塞莱斯提乌斯被佐西姆谴责。}

可敬的教宗佐西姆，考虑到这一恳求性的前言，处理了这位被错误学说的气息所膨胀的人，使他应该谴责所有由执事保利努斯对他提出的反对点，并应赞同由他前任（神圣记忆）发出的宗座诏书。然而，被告拒绝谴责执事提出的反对，但他不敢抵抗真福教宗依诺增爵的信；事实上，他甚至\Quote{承诺他将谴责宗座所谴责的所有点}。因此，这人被温和的治疗对待，像一个需要休息的谵妄病人；但同时，他还不被认为已准备好从绝罚的束缚中释放。给予两个月的间隔，直到能收到来自非洲的通信，一个恢复的地方被让与他，在已宣布的判决的温和复原下。因为事实上，如果他愿意放下他虚妄的固执，现在愿意执行他所承担的事，并仔细阅读他通过承诺服从所回复的那封信，他本会变得更好。但在非洲主教会议正式发出诏书之后，有充分的理由按照最严格的公平对他执行判决。这些理由是什么，你们可以自己阅读，因为我们已将一切详情发送给你们。\par

% structure: p0019 source=h2 target=subsection reason=relative-heading-rank; base=h2

\subsection{第九章 [VIII.]——佩拉纠欺骗了巴勒斯坦的会议，但未能欺骗罗马教会。}

因此，佩拉吉乌斯若能真诚反思自己的立场与著作，也无理由说他不该受到这样的判决。因为他虽然欺骗了巴勒斯坦的会议，表面上在那里为自己辩护，却完全未能瞒过罗马的教会（如您所知，他在那里绝非陌生人），尽管他尽力尝试，或许能得逞。但正如我刚才所说，他完全失败了。因为最蒙福的教宗佐西默记得他的前任——那位为他树立了如此值得效仿的榜样——对同一事件的意见。他也没有忽视忠信的罗马人对这人的看法，他们的信德在主内值得称扬\BibleRef{罗 1:8}，并且他看到他们对捍卫公教真理反对这一异端的热忱，既热烈又极其一致。这人在他们中间生活已久，他的观点不可能逃过他们的注意；此外，他们已完全识破了他的门徒塞莱斯提乌斯，以至于能够立即提出最可信且不可反驳的证据。至于圣教宗依诺增爵对巴勒斯坦会议（佩拉吉乌斯自夸已被该会议宣告无罪）的庄严判决，您可以在他写给我们的信中读到。非洲会议给可敬的教宗佐西默的回信中也适当提到了此事，我们已随其他指示寄给了您们。但与此同时，我认为不应省略在此作品中引出具体内容。\par

% structure: p0021 source=h2 target=subsection reason=relative-heading-rank; base=h2

\subsection{第10章[IX.]——依诺增爵对巴勒斯坦会议的判决}

于是，五位主教（我是其中之一）给他写了一封信，其中我们提到了巴勒斯坦的会议，该会议的报告已传到我们耳中。我们告诉他，在这人居住的东方，发生了一些教会会议程序，他在这些程序中被认为已就所有指控被宣告无罪。依诺增爵对我们这封函件的回信中有以下的话：\Quote{在这些会议记录本身中，有若干论点，当反对他时，他部分地通过回避加以压制，部分地通过曲解许多词语的意义而使其完全模糊；还有其他指控他——不是以他当时似乎使用的诚实方式，而是以诡辩的方法——予以澄清，对某些反对意见断然否认，对另一些则以错误的解释加以篡改。然而，但愿他现在甚至采取更为可取的途径，从自己的错误转回公教信仰的真道；但愿他也能妥善考虑天主的日常恩宠，承认恩宠的帮助，愿意并渴望在全人的赞同下，通过公开信服的方法——不是通过司法程序，而是通过真心皈依公教信仰——真正得到纠正。因此，我们既不能赞同也不能责备该审判中的程序；因为我们不知道程序是否真实，即便真实，它们是否确实表明这人通过逃避而非通过完全的真理为自己辩护。}您从这些话中可以清楚看出，最蒙福的教宗依诺增爵无疑将此视为一个对他来说绝非陌生的人。您可以看到他对这人的无罪宣判持何种意见。您还可以看到，他的继任者圣教宗佐西默必须记住什么——正如他确实记住的——从而毫不犹豫地确认其前任在此案中的判决。\par

% structure: p0023 source=h2 target=subsection reason=relative-heading-rank; base=h2

\subsection{第11章[X.]——佩拉吉乌斯如何欺骗巴勒斯坦会议}

现在我恳请您仔细注意，佩拉吉乌斯被证明是如何欺骗巴勒斯坦的审判官的——且不提其他要点，仅就婴儿圣洗这一问题，以免有人以为我们使用了诽谤和怀疑，而非确知的事实，当我们断言佩拉吉乌斯隐瞒了塞莱斯提乌斯更坦率表达的意见，而实际上他怀有相同的观点。从上述内容已清楚看到，塞莱斯提乌斯拒绝谴责\Quote{亚当的罪只伤害他自己，而不伤害人类，并且婴儿出生时的状态与亚当犯罪前的状态相同}这一断言，因为他明白，如果他谴责这些命题，他就得承认婴儿身上有来自亚当的罪的遗传。然而，当有人反对佩拉吉乌斯，说他在这一点上与塞莱斯提乌斯意见一致时，他毫不犹豫地谴责了那些话。我很清楚您以前读过这一切。但既然我们不只是为您一个人写这篇叙述，我们继续转录会议记录本身的原文，免得读者不愿自己查阅记录，或者如果没有记录，不愿费心获取副本。以下是原文：——\par

% structure: p0025 source=h2 target=subsection reason=relative-heading-rank; base=h2

\subsection{第12章[XI.]——巴勒斯坦会议关于佩拉吉乌斯案件的记录摘录}

\Quote{会议说：如今，因为\Person{白拉奇}已就这一不确定的愚昧之言宣布绝罚，正确地回答说，人藉天主的帮助和恩宠能过\OriginalTerm{无罪}{\GreekText{ἀναμάρτητος}}的生活，也就是说，成为无罪的人，就请他也在其他条款上给我们答复。\Person{塞勒斯提乌}（\Person{白拉奇}的门徒）的教导中另一个细节，选自\Place{迦太基}的圣\Person{奥勒留}主教及其他主教面前所提及并听过的条目，其内容如下：\InnerQuote{亚当被造成有死的，并且无论他犯罪与否都会死；亚当的罪只伤害他本人，不伤害人类；法律与福音一样引导我们进入天国；在基督来临之前已有无罪的人；新生婴儿与亚当在犯罪前的状况相同；一方面，整个人类并不因亚当的死亡和过犯而死，另一方面，整个人类也不因基督的复活而复活；圣\Person{奥斯定}主教写了一本书回答他在\Place{西西里}的追随者，针对所附的条款，而在这本写给\Person{希拉略}的书中包含了以下陈述：一个人若愿意就能无罪；婴儿即使未受洗也有永生；富有的基督徒即使受了洗，除非他们放弃并舍弃一切，否则无论他们似乎行了什么善事，都不得算为他们所有，也不能拥有天国。}然后\Person{白拉奇}说：关于人能否无罪，我的意见已经说过了。然而，关于所声称的主来临之前已有过无罪生活的人，我们同时也说，在基督来临之前，确实有人过圣洁正义的生活，正如圣经所传给我们的记载。至于其他各点，即使按他们自己的说法，其性质亦无需我负责；但为满足神圣会议起见，我宣布绝罚那些目前或曾经持有这些意见的人。}\par

% structure: p0027 source=h2 target=subsection reason=relative-heading-rank; base=h2

\subsection{第十三章——\Person{塞勒斯提乌}是更胆大的异端者；\Person{白拉奇}是更狡猾的。}

你确实看到，且不提其他各点，\Person{白拉奇}如何宣布绝罚那些认为\Quote{亚当的罪只伤害他本人，而不伤害人类；并且婴儿出生时的状况与亚当犯罪前的状况相同}的人。那么，审判他的主教们还能理解为他是什么意思呢？难道不是理解为亚当的罪传给了婴儿吗？\Person{塞勒斯提乌}正是为了避免作这样的承认，才拒绝谴责这一陈述，而此人（指\Person{白拉奇}）却相反地绝罚了它。因此，如果我要证明，他（\Person{白拉奇}）对婴儿的看法实际上与\Person{塞勒斯提乌}没有什么两样，即认为婴儿出生时没有任何罪的沾染，那么，在这个问题上，他和\Person{塞勒斯提乌}之间还有什么区别呢？无非是，一个更公开，一个更含蓄；一个更固执，一个更虚伪；或者至少可以说，一个更坦率，一个更狡猾。因为，前者（\Person{塞勒斯提乌}）在\Place{迦太基}教会面前拒绝谴责，后来却在\Place{罗马}教会公开承认那是他自己的主张；同时，他自称\Quote{如果错误不知不觉地进入了他，他愿意接受纠正，因为他不过是人}；而后者（\Person{白拉奇}），虽然谴责了这一教条，认为它与真理相悖，免得被他的公教审判者们谴责，但却保留以备日后辩护，因此，或者他的谴责是谎言，或者他的解释是诡计。\par

% structure: p0029 source=h2 target=subsection reason=relative-heading-rank; base=h2

\subsection{第十四章——他表明，即使在\Place{巴勒斯坦}会议之后，\Person{白拉奇}在原罪问题上仍持有与\Person{塞勒斯提乌}相同的见解。}

然而，我看得出，人们最有理由要求我不可拖延我所作的承诺，即证明他实际上怀有与\Person{塞勒斯提乌斯}相同的观点。在他较近期的作品（为辩护自由意志而写）的第一卷中（他在寄往\Place{罗马}的信中提到了这部作品），他说：\Quote{一切善与一切恶，凡使我们因此值得赞美或责备的，都不是与生俱来的，而是我们做出来的；因为我们出生时并非已完全发展，而是具有行善或行恶的能力；我们被生育出来时既无德行，亦无恶习；在我们自身意志的行动之前，人身上只有天主所形成的一切。}现在您看出，在\Person{佩拉吉乌斯}的这些话里，包含了这两个人的教义：即婴儿出生时没有沾染\Person{亚当}的任何罪。因此，\Person{塞勒斯提乌斯}拒绝谴责那些说\Person{亚当}的罪只伤害了他自己、而非整个人类的人，以及那些说婴儿出生时的状态与\Person{亚当}在违命之前的状态相同的人，这并不奇怪。但非常令人惊讶的是，\Person{佩拉吉乌斯}竟厚颜无耻地诅咒这些意见。因为，如果如他所说，\Quote{恶不是与我们同生的，我们被生育出来时没有过犯，人身上在自己意志行动之前只有天主所形成的一切，}那么当然\Person{亚当}的罪只伤害了他自己，因为它没有传给后代。因为没有一种罪不是恶，或者一种罪不是过犯；否则罪就是天主所造的。但他说：\Quote{恶不是与我们同生的，我们被生育出来时没有过犯，人在出生时只有天主所形成的一切。}既然他用这话认为他那句众所周知的话\Quote{亚当的罪只伤害了他自己、而非整个人类}是千真万确的，那么为什么\Person{佩拉吉乌斯}要谴责这一点呢？若不是为了欺骗他的公教法官们？同理，也可以论证：\Quote{如果恶不是与我们同生的，如果我们被生育出来时没有过犯，如果人在出生时只有天主所形成的一切，}那么毫无疑问\Quote{婴儿出生时的状态与亚当在违命之前的状态相同}，在他身上没有固有的恶或过犯，只有天主所形成的一切。然而\Person{佩拉吉乌斯}却诅咒一切\Quote{现在或过去任何时候认为新生的婴儿出生时的状态与亚当在违命之前的状态相同}的人——换言之，即没有任何恶，没有任何过犯，只有天主所形成的一切。那么，\Person{佩拉吉乌斯}为什么又谴责这个论点呢？若不是为了欺骗公教会会议，并使自己免于被定为异端新说者的罪？\par

% structure: p0031 source=h2 target=subsection reason=relative-heading-rank; base=h2

\subsection{第十五章——\Person{佩拉吉乌斯}以其谎言和欺骗从\Place{巴勒斯坦}主教会议窃取无罪开释。}

至于我本人，正如您所熟知的，并且在我写给可敬的、年迈的\Person{奥勒留斯}关于\Place{巴勒斯坦}会议进程的书中也已说明的，我确实感到高兴，\Person{佩拉吉乌斯}在那次答辩中已经穷尽了这个问题。在我看来，他通过宣布诅咒那些认为因\Person{亚当}的罪只有他自己、而非整个人类受到伤害的人，以及那些认为婴儿与第一人在违命之前处于相同状态的人，似乎已经最清楚地承认了婴儿有原罪。然而，当我读了他的四卷书（我从第一卷中抄录了我刚才引用的那些话），并发现他仍然怀有反对公教信仰关于婴儿的见解时，我对他在教会的会议上、在如此重大的问题上竟如此无耻地坚持谎言感到更加惊讶。因为如果他早已写了这些书，他怎能声称诅咒那些曾经持有上述意见的人呢？然而，如果他打算以后出版这样的作品，他又怎能诅咒那些当时持有这些意见的人呢？除非，他凭某种可笑的遁词，意指他所诅咒的对象是那些在过去某个时候曾经持有、或当时正在持有这些意见的人；但至于将来——即那些将要接受这类观点的人——他觉得他不可能对自己或他人预作判断，因此当后来被发现坚持类似错误时，他并没有撒谎。不过，他并没有提出这个借口，不仅因为它可笑，而且因为它不可能是真的；因为正是在这些书中，他既论证罪不从\Person{亚当}传给婴儿，又夸耀\Place{巴勒斯坦}主教会议的进程，在那里他被认为真诚地诅咒了持有争议意见的人，而事实上，他通过施行欺骗窃取了他的无罪开释。\par

% structure: p0033 source=h2 target=subsection reason=relative-heading-rank; base=h2

\subsection{第十六章——\Person{佩拉吉乌斯}欺诈狡猾的借口。}

因为当我们现在所要处理的问题中，他回答追随者的话究竟有什么意义呢？他告诉他们：\Quote{他之所以谴责那些反对他的论点，是因为他自己坚持认为，原罪不仅对第一个人有害，而且对整个人类有害，不是通过遗传，而是通过榜样；}换句话说，并非所有从他繁衍而来的人都从他那里继承了任何过犯，而是所有后来犯罪的人都效法了犯第一个罪的人？或者当他说：\Quote{婴儿之所以不在亚当堕落前的状态中，是因为他们还不能领受诫命，而亚当能；还因为他们尚未运用理性意志的选择，而亚当确实运用了，否则就不会给他诫命}？这样的解释如何使他觉得有理由正确谴责这些命题：\Quote{亚当的罪只伤害了他自己，而没有伤害全人类}和\Quote{婴儿出生时与亚当犯罪前的状态完全相同}？并且通过上述谴责，他在后来著作中持有这些观点时并非有欺骗之罪？即\Quote{婴儿出生时没有任何邪恶或过犯，他们里面只有天主所形成的东西}——简言之，没有敌人造成的创伤？\par

% structure: p0035 source=h2 target=subsection reason=relative-heading-rank; base=h2

\subsection{第十七章——佩拉基乌斯如何欺骗了他的审判官}

现在，他是不是通过这样陈述，以另一种意义的解释来反驳被提出的反对意见，从而证明他没有欺骗那些审判他的人？那么他完全达不到目的。他的解释越狡猾，他欺骗他们的手段就越隐蔽。因为，正因为他们是公教主教，当他们听到那人绝罚那些主张\Quote{亚当的罪只伤害了他自己，而没有伤害全人类}的人时，他们理解为他只是宣告公教会惯常所教导的，即教会确实为婴儿施洗以赦免罪过——不是因模仿第一个罪人的榜样而犯的罪，而是因生育本身、因起源的腐化而沾染的罪。当他们又听到他绝罚那些主张\Quote{婴儿出生时与亚当堕落前的状态相同}的人时，他们以为他指的是那些\Quote{认为婴儿没有从亚当继承任何罪，因此他们处于亚当犯罪前的状态}的人。当然，除了论点所围绕的那个异议之外，不会有其他反对他的异议。因此，当他这样解释异议，说婴儿之所以不在亚当犯罪前的状态，仅仅是因为他们还没有达到同样的身心坚固程度，而不是因为任何遗传的过犯传递给他们时，必须这样回答他：\Quote{当这些异议被提交给你要求谴责时，公教主教们并未以此意义理解它们；因此，当你谴责它们时，他们相信你是公教徒。因此，他们认为你所持的立场应免于责备；但你实际所持的立场应受谴责。所以，被宣告无罪的不是你，因为你持有应受谴责的教义；而是你本应持有并坚持的意见被免除了责备。你只能被认为是被宣告无罪，因为你被认为持有值得称赞的见解；因为你的审判官无法想像你隐瞒了应受谴责的见解。你已被公正地裁定为塞勒斯提乌斯的同谋，因为你证明自己同享他的观点。虽然你在审判期间封闭了你的书籍，但在审判后你却将它们公之于世。}\par

% structure: p0037 source=h2 target=subsection reason=relative-heading-rank; base=h2

\subsection{第十八章——对佩拉基乌斯的定罪}

既然如此，你当然会感到，主教会议、宗座、整个罗马教会，以及因天主恩宠而成为基督徒的罗马帝国本身，都最正义地起来反对这一邪恶错误的创始人，直到他们悔改并逃脱魔鬼的陷阱。因为谁能知道天主会不会赐给他们悔改的心，使他们发现、承认甚至宣扬祂的真理\BibleRef{弟后 2:25}，并谴责他们自己该受诅咒的错误？但无论他们自己的意志如何，我们毫不怀疑，主的仁慈眷顾已经为许多跟随他们的人谋取了益处，唯一的原因是看到他们与公教会共融。\par

% structure: p0039 source=h2 target=subsection reason=relative-heading-rank; base=h2

\subsection{第十九章——佩拉基乌斯企图欺骗宗座；他颠倒了争论的关键点}

但我愿你们仔细留意\Person{白拉奇}如何企图用欺骗的手段，甚至在这婴儿圣洗的问题上，瞒过宗徒座的主教的判断。他将一封信寄往罗马，给有福的纪念的教宗\Person{依诺增爵}；当信到达时，教宗已不在人世，于是信被交到圣教宗\Person{佐西摩}手中，并由他转交给我们。在这封信中，他抱怨说，\Quote{某些人诽谤他，说他不给婴儿施行圣洗圣事，并许诺无论基督的救赎如何，都能得到天国。}然而，反对他们的人并非如他所说的那样提出这些指责。因为他们既不否认给婴儿施行圣洗圣事，也不许诺任何人可以脱离基督的救赎而得到天国。因此，关于他抱怨被某些人诽谤一事，他如此陈述，以便能够轻易地回答对他的指控，而不损害自己的教义。[第十八章]真正反对他们的是，他们拒绝承认未受圣洗的婴儿应受第一个人的定罪，并且原罪已传给他们，需要藉重生来洗净；他们主张婴儿受圣洗仅仅是为了进入天国，好像只有那些不能领受主身体和血的人，才会在远离天国时遭遇永死。我愿你们知道，这就是关于婴儿圣洗的真正反对意见；而不是像他所陈述的那样，目的是为了使自己能够保全自己的教义，同时以回应反对意见为借口，回答实际上是他自己提出的主张。\par

% structure: p0041 source=h2 target=subsection reason=relative-heading-rank; base=h2

\subsection{第二十章——\Person{白拉奇}在含混的遁词中为他的谬误提供庇护}

接着，请看他是如何回答的，如何在他双重意义的晦涩迷宫中为他的虚假教义提供避难所，用错误的黑暗迷雾扑灭真理；以致就连我们，在第一次读他的话语时，也几乎为它们的恰当和正确而欣喜。但他书籍中更充分的讨论（在这些讨论中，他通常尽管竭力隐藏，却被迫解释自己的意思）使我们甚至对他那些较好的陈述也产生怀疑，生怕在仔细审视时发现它们是含混的。因为，在说了\Quote{他从未听过甚至一个不敬神的异端者说过这话}（即他作为反对意见所提出的关于婴儿的话）之后，他接着问道：\Quote{谁对福音教训如此生疏，不仅试图作这样的肯定，甚至能轻易地说出它，或让它进入他的思想？然后，谁如此不敬神，竟想禁止婴儿受圣洗并在基督内重生，以排除他们于天国之外？}\par

% structure: p0043 source=h2 target=subsection reason=relative-heading-rank; base=h2

\subsection{第二十一章[第十九]——\Person{白拉奇}回避了婴儿为何需要\OriginalTerm{圣洗}{Baptism}的问题}

他说这一切毫无用处。他并没有因此为自己辩白。就连他们，也从未否认婴儿没有圣洗不能进入天国。但这并不是问题所在；我们所讨论的是婴儿身上原罪的清除。让他就此点为自己辩白吧，因为他拒绝承认婴儿身上有什么需要重生之水洗涤的东西。因此，我们应当仔细考虑他随后要说的。在引用了福音中宣布\BibleQuote{凡不由水和圣神重生者，不能进天国}\BibleRef{若 3:5}的经文之后（关于这一点，正如我们所说，他们不提出质疑），他立刻问道：\Quote{谁如此不敬神，竟忍心拒绝任何年龄的婴儿分享人类的共同\OriginalTerm{救赎}{redemption}？}但这是含混的言语；因为他所说的\Quote{救赎}是什么意思？是从恶到善？还是从善到更好？因为连\Person{塞勒斯提乌斯}在迦太基，也在他的书中允许婴儿有救赎；尽管同时，他不承认从亚当传来的罪到他们身上。\par

% structure: p0045 source=h2 target=subsection reason=relative-heading-rank; base=h2

\subsection{第二十二章[第二十]——\Person{白拉奇}含混的另一个实例}

接着，再看他在最后一句评论之下添加的话：\Quote{有谁，}他说，\Quote{能禁止那已生于这\OriginalTerm{不确定的生命}{uncertain life}的人，重生而进入永恒而确定的生命？}换言之：\Quote{谁如此不敬神，竟禁止那已生于这不确定生命的人，重生而进入那确定而永恒的生命？}当我们第一次读这些话时，我们以为他用\Quote{不确定的生命}一词意指这现世暂时的生命；尽管在我们看来，他更应当称它为\Quote{可朽的}而不是\Quote{不确定的}，因为它是被确定的死亡所终结的。但尽管如此，我们仍然认为他更喜欢称这可朽的生命为\Quote{不确定的}，因为人们普遍认为，在我们的生命中确实没有一刻是免于这种不确定性的。因此，我们的焦虑在某种程度上被以下想法缓解了，这个想法几乎成了证据：尽管他不愿公开承认未领受圣洗圣事而离世的婴儿会遭受\OriginalTerm{永死}{eternal death}，但我们论证说：\Quote{如果，正如他似乎承认的，只有受圣洗的人才能获得永生，那么自然，未受圣洗而死的人就陷入永死。}然而，这种命运决不能公正地临到那些在今生从未犯过任何自身之罪的人，除非是由于原罪。\par

% structure: p0047 source=h2 target=subsection reason=relative-heading-rank; base=h2

\subsection{第二十三章[第二十一]——他所谓的我们出生进入的\Quote{不确定的生命}是什么意思？}

然而，有些弟兄后来并未忘记提醒我们，伯拉纠之所以可能这样表达，是因为他据说在此问题上对一切询问者都有现成的回答，即：\Quote{至于那些未受圣洗即死去的婴儿，我确实知道他们不去哪里；但至于他们去哪里，我不知道。}也就是说，他知道他们不去天国。但至于他们去哪里，他过去（而且就这事而言，\Emph{现在仍}）习惯说不知道，因为他不敢说那些他深信在这世上从未犯过罪、且他不承认继承原罪的人会去永死。因此，他那些被送到罗马以确保他完全无罪的开脱之词，充满了模棱两可，为其教义提供了庇护，从中可跳出一种异端的意义，诱捕粗心的迷途者；因为当无人能给出答案时，任何独处者都可能发现自己软弱。\par

% structure: p0049 source=h2 target=subsection reason=relative-heading-rank; base=h2

\subsection{第24章——伯拉纠长期居住在罗马}

事实上，在他写给上述教宗依诺增爵的信中，连同这封信一同送到罗马的信仰之书里，他不过是更明显地在掩饰中暴露了自己。他说：\Quote{我们持守同一圣洗，我们认为在婴儿身上和成人身上应以同样的圣事言语施行。}然而，他并没有说\Quote{同一圣事}（即便他如此说，也仍会有歧义），而是\Quote{同一圣事言语}——仿佛婴儿罪过的赦免是由言语的声音宣告，而非因礼节的效用实现。当时他看似说出了符合大公信仰的话；但他未能永久欺骗那宗座。在非洲公会议（该教省已悄然蔓延了这种有害教义——但未广为流布或深入扎根）的复函之后，此人的其他见解也被一些忠信弟兄的勤勉发现并带到罗马，他曾在罗马居住很长时间，并已参与过各种论述和争论。为使这些见解遭到谴责，教宗佐西莫斯（如你所读）将之附在他为在全大公世界颁布而写的信中。在这些陈述中，伯拉纠假托讲解宗徒保禄的\BibleBook{}{致罗马人书}，用这些言辞论证：\Quote{如果亚当的罪伤害了那些没有犯罪的人，那么基督的义也裨益那些不信仰的人。}他也说了其他类似的话；但我在所著\WorkTitle{论婴儿圣洗}一书中，已靠主的帮助全部驳斥并回答了它们。然而在之前所谓的注释中，他不敢以自己的名义说出这些可指责的陈述。但这一特定陈述在如此熟悉他的地方提出后，他的言语和意思就无法掩饰了。在那些书（我在前面已引过其第一部）中，他毫不掩饰地论述了这一点。他竭尽全力最明确地断言，婴儿的人性绝不能设想为因生育而败坏；并且他为他们要求救恩作为应得之物，从而藐视了救主。\par

% structure: p0051 source=h2 target=subsection reason=relative-heading-rank; base=h2

\subsection{第25章——伯拉纠和塞勒斯提乌斯的谴责}

既然如此，我陈述的这些事现在清楚表明，出现了致命的异端，靠着主的帮助，教会此时更为直接地防范它——如今那两人伯拉纠和塞勒斯提乌斯，或是已被给予补赎的机会，或是因拒绝而完全被谴责。据传，或者可能已被证实，他们是这种偏离的作者；无论如何，若非作者（因从他人学得），他们仍是其自我夸耀的拥护者和导师，通过他们，异端已发展并传播得更广。这种夸耀甚至体现在他们自己的陈述和著作中，以及现实无可置疑的标志中，还有由此产生并增长的名声中。那么，还剩下什么？难道不是每位大公信徒都当以主所赐予的一切力量驳斥这种有害教义，并以一切警惕反对它；这样，每当我们为真理争辩（被迫回答，但并不喜爱争论），无知者可得教导，教会也因敌人为她毁灭而设的计谋获益，正如宗徒的话所说的：\Quote{必须有异端，好使那些经得起考验的人得以显明在你们中间}？\BibleRef{格前 11:19}\par

% structure: p0053 source=h2 target=subsection reason=relative-heading-rank; base=h2

\subsection{第26章——伯拉纠派认为提出关于原罪的问题不会危及信仰}

因此，在我们以书面形式充分讨论了如何反驳他们这种极其敌对天主藉着我们的主耶稣基督赐给大小之人的恩宠的错误之后，现在我们有责任审视并驳斥他们巧妙提出的那个断言，即\Quote{将这个问题提出来讨论不会对信德造成危险}——目的是为了使他们看起来，如果被发现偏离了信德，则似乎不是犯了有罪的错误，而只是，可以这么说，礼貌上的失误。这正是塞勒斯提乌斯在迦太基的教会诉讼程序中所说的话：\Quote{关于罪的传袭，} 他说，\Quote{我已经说过，我听过许多在天主教会中公认有地位的人否认它，而另一方面也有许多人肯定它；它确实可以公平地视为一个可探究的问题，而不是异端。我一直坚持婴儿需要圣洗，并应受洗。他还想要什么？} 他这么说，似乎是想暗示，只有当他说婴儿不应受洗时，他才会被视为可指控为异端。然而就情形而言，既然他承认他们应受洗，他认为自己没有（有罪地）犯错，因此不应被判为异端者，即使他坚持关于婴儿受洗的理由与真理所持或信德所要求的不同。按照同样的原则，在他送到罗马的那本书中，他首先按照自己的喜好解释了自己的信仰，从独一天主性的天主圣三直到将来死者的复活方式；然而，关于所有这些点，从未有人问过他，他也没有问过别人。当他的论述达到此刻正在考虑的问题时，他说：\Quote{如果确实有任何超过信德范围的问题出现，可能有许多人对此有分歧，我绝没有假装对任何教义作出判决，仿佛我自己拥有在这件事上的最终权威；但我从先知和宗徒的泉源中得来的一切，都呈交给您的宗座予以核准，以便若有任何错误因着我们的无知而潜入我们中间（因为我们是人），可以由您的裁决来纠正。} 你当然清楚地看到，在他的这一行动中，他使用了所有这些恳求性的前言，以便如果他被发现有任何错误，他可能显得不是在信德问题上犯错，而是在信德之外的可疑点上犯错；其中，无论纠正错误多么必要，但不是作为异端来纠正；同样，受纠正的人确实被宣告有错误，但仍然不被判为异端者。\par

% structure: p0055 source=h2 target=subsection reason=relative-heading-rank; base=h2

\subsection{第27章 [XXIII.]——论信德之外的问题：它们是什么及其实例。}

但他在这个观点上大错特错。他所认为的信德之外的问题，与那些在不损害我们成为基督徒的信德的情况下，存在着对事实的不知因而对任何确定意见有所保留，或者由于人类思想的软弱而产生不符合真理的猜测性观点的情形截然不同。例如：当有人质疑天主从泥土中造人并放置于其中的那个乐园的描述和位置时——尽管这丝毫不会动摇基督徒信仰相信确实存在这样一个乐园；或者当有人问厄里亚和哈诺客此刻在哪里——是在这个乐园中还是别处，尽管我们毫不怀疑他们仍以出生时的同一身体存在；或者当有人询问宗徒被提到第三层天上时是在身内还是在身外——然而，这种询问对那些人来说极不谦虚，他们想知道正是这奥秘本身的主体明确宣称不知道的事\BibleRef{格后 12:2}，但并不损害他对事实的信仰；或者当有人提出疑问：他告诉我们他被提到的那\Quote{第三}层天究竟有多少层；或者这个可见世界的元素是四个还是更多；是什么导致太阳或月亮发生天文学家习惯预测的特定季节的食相；为什么古代人能活到圣经所记载的年龄；他们生育头生子达到青春期的时期是否因寿命更长而推迟到更年长的年龄；默突舍拉可能住在哪里，因为他不在方舟里——根据大多数圣经抄本（希腊文和拉丁文）的年代记录，他被发现活到了洪水之后；或者我们是否必须遵循少数抄本（它们恰好极少）的顺序，这些抄本安排年份显示他在洪水之前就死了。在种种此类纷繁无数的问题中，涉及天主最隐秘的作为或圣经中最隐晦的段落，难以用任何确定的方式把握和界定，谁不感到：在许多点上，无知可以与纯正基督徒信德共存，有时也可以持有错误意见而不至被指控为异端道理？\par

% structure: p0057 source=h2 target=subsection reason=relative-heading-rank; base=h2

\subsection{第28章 [XXIV.]——\Person{白拉奇}和\Person{塞勒斯提乌斯}的异端攻击我们信德的根基。}

然而，就两个人而言，我们因其中一人被卖于罪下\BibleRef{罗 7:14}，因另一人从罪中被救赎——因前者我们堕入死亡，因后者我们获得生命；前者在自身中毁灭了我们，因为他行自己的意愿，而非造物主的意愿；后者在自身中拯救了我们，因为祂不行自己的意愿，而行了派遣祂者的意愿。基督徒信仰确切地在于这两个人。因为\BibleQuote{只有一个天主，在天主与人之间只有一个中保，就是降生成人的基督耶稣}\BibleRef{弟前 2:5}；因为\BibleQuote{除他以外，无论凭谁，决无救援，因为在天下人间，没有赐下别的名字，使我们赖以得救的}\BibleRef{宗 4:12}；并且\BibleQuote{天主在祂内为众人确定了信仰，因为祂使祂从死者中复活}\BibleRef{宗 17:31}。没有这信德——也就是说，不信在人与天主之间的唯一中保、降生成人的基督耶稣；我再说，不信祂的复活，天主借此给众人确证，而若没有祂的降生成人和死亡，无人能真正相信；因此，若不信基督的降生成人、死亡和复活，则基督信仰真理毫不迟疑地宣告：古圣们不可能从罪中被洁净而成为圣洁，并因天主的恩宠而成义。这既适用于圣经中提到的圣人，也适用于那些虽未被提及但必然存在的——或是在洪水之前，或是在洪水与赐予法律之间，或是在法律本身时期——不仅包括以色列子民中的先知，甚至包括那民族之外的人，例如约伯。因为这些人的心也是因对同一中保的信德被洁净，并且\BibleQuote{天主的爱藉着所赐与他们的圣神，已倾注在他们心中}\BibleRef{罗 5:5}，\BibleQuote{圣神随意向那里吹}\BibleRef{若 3:8}，并不随从人的功德，而是亲自产生这些功德。因为天主的恩宠除非完全自由，否则便根本不存在。\par

% structure: p0059 source=h2 target=subsection reason=relative-heading-rank; base=h2

\subsection{生活在法律时期的义人并非在法律之下，而是在恩宠之下。新约的恩宠隐藏于旧约之下。}

死亡从亚当直到梅瑟确实掌权了\BibleRef{罗 5:14}，因为甚至通过梅瑟所赐的法律也不能战胜它：事实上，法律并非赐下作为能赐予生命的东西\BibleRef{迦 3:21}，而是应当向那些死去的人、那些需要恩宠赐予生命的人表明，他们不仅被罪的繁衍和统治所压倒，而且还因违反法律本身的附加罪过而被定罪：这并非为了使任何人在天主的怜悯中（即使在那早期）理解这事而丧亡；而是为了使那注定因死亡的统治而受罚、又因自己违反法律而显为有罪的人，能寻求天主的帮助，这样，罪恶在那里增多，恩宠在那里便更加增多\BibleRef{罗 5:20}，即那唯独从这死亡的身体中拯救的恩宠\BibleRef{罗 7:24}。[XXV.]然而，尽管如此，虽然甚至梅瑟所赐的法律也不能将任何人从死亡的统治中解放出来，但在法律时期，那时也有天主的子民，他们不是生活在法律的惊惧、定罪和惩罚之下，而是生活在恩宠的喜悦、医治和解放之中。有些人说过：\BibleQuote{我是在罪孽中成形的，我母亲在罪中怀了我}；以及\BibleQuote{因我的罪，我骨中不得安宁}；以及\BibleQuote{天主啊，求祢在我内里造一颗清洁的心，使我灵里重新有正直的灵}；以及\BibleQuote{以祢引导的圣神坚固我}；以及\BibleQuote{不要从祢收回祢的圣神}。还有一些人说过：\BibleQuote{我信了，所以我说了话}。因为他们也以我们所用的同一信德得以洁净。因此宗徒也说：\BibleQuote{我们既然有同一信德的精神，正如经上记载：\InnerQuote{我信了，所以我说了话}，我们也信，所以也说}\BibleRef{格后 4:13}。正是出于同一信德，才说：\BibleQuote{看，一位贞女将怀孕生子，人要称祂的名字为厄玛奴耳}\BibleRef{依 7:14}，\Quote{解释出来就是：天主与我们同在}\BibleRef{玛 1:23}。也正是出于信德，论及祂说：\BibleQuote{祂像新郎走出洞房，像巨人欢然奔跑祂的路程。祂的出发从天的尽头，祂的回转直达天的另一端；没有一个人能躲过祂的热力}。又是出于信德，向祂说：\BibleQuote{天主啊，祢的宝座永远常存，祢王国的权杖是正直的权杖。祢喜爱正义，憎恨不义；因此天主，祢的天主，用喜乐之油傅了祢，超过祢的同伴}。藉着同一信德之神，他们预见了这一切事将要发生，而我们则相信它们已经发生了。确实，那些能以忠信之爱向我们预言这些事的人，自己并没有分沾它们。宗徒伯多禄说：\BibleQuote{你们为什么试探天主，要把一个轭放在门徒的颈上，这轭连我们的祖先和我们都不能承担？但我们相信，我们得救是因主耶稣基督的恩宠，和他们一样}\BibleRef{宗 15:10}。他说这话的根据是什么，难道不是因为就连他们也是因主耶稣基督的恩宠得救，而非因梅瑟的法律？法律带来的不是治愈，而只是对罪的认识\BibleRef{罗 3:20}。但如今，天主的正义在法律之外已显明出来，被法律和先知所证明\BibleRef{罗 3:21}。因此，如果它现已显明，那么它那时就已存在，只是隐藏着。这隐藏是由圣殿的帐幔所象征。基督死亡时，这帐幔裂为两半\BibleRef{玛 27:51}，以表明祂的完全启示。因此，从古时起，在天主的子民中就已存在着这唯一中保——天主与人之间的中保，即是人基督耶稣——的恩宠；但如同雨落在羊毛上——那是天主为自己产业所分别出来的，不是出于债务，而是出于祂自己的旨意——它曾是隐秘地临在，如今却显而易见地显于万民之中，如同它的\Quote{禾场}，而羊毛是干的——换句话说，犹太民族成了被弃绝的\BibleRef{民 6:36}。\par

% structure: p0061 source=h2 target=subsection reason=relative-heading-rank; base=h2

\subsection{第三十章 [XXVI]——佩拉吉乌斯和塞莱斯提乌斯否认古代圣人是因基督得救}

因此，我们不可像佩拉吉乌斯及其门徒那样划分时代，他们说：人首先按本性过义的生活，然后是在法律之下，第三是在恩宠之下——本性是指从亚当到法律颁赐之前的全部漫长时期。\Quote{那时，}他们说，\Quote{造物主藉理性的引导被人认识；生活正直的规则写在人的心中，不是写在文字的法律中，而是写在自然的法律中。但人的品行变得败坏；然后，}他们说，\Quote{当本性如今已玷污，开始不足时，法律被加给它，如同藉着一轮月亮，在本性的红晕受损之后，恢复了它原有的光泽。但当犯罪的习惯在人间太过盛行，法律不足以治愈它时，基督来了；这位医生亲自藉着祂自己，而非藉着祂的门徒，在那疾病最危急的发展阶段带来了缓解。}\par

% structure: p0063 source=h2 target=subsection reason=relative-heading-rank; base=h2

\subsection{第三十一章——基督的降生成人对于先祖们是有益的，即使它尚未发生}

通过这种争辩，他们试图将古圣祖排除在中保的恩宠之外，仿佛人基督耶稣不是天主与\Emph{那些人}之间的中保；理由是在那些义人生活的时候，祂尚未从童贞女的子宫中取得血肉，还不是人。然而，如果这是真的，那么宗徒所说的就徒然了：\BibleQuote{死因一人而来，复活也因一人而来；因为正如在亚当内众人都死了，照样在基督内众人也都要复活。}\BibleRef{格前15:21}因为按照这些人的虚妄想法，那些古圣祖认为他们的本性是自足的，不需要人基督作他们的中保来与天主和好，所以他们也不会在祂内复活，因为他们显然不属于祂的身体，作为肢体，按照那说法：祂成为人是为了人。然而，正如真理通过祂的宗徒所说，既然在亚当内众人都死了，照样在基督内众人也都要复活；因为复活来自一个人，正如死亡来自另一个人；有哪个基督徒敢怀疑：即使在人类更远古时期中悦天主的那些义人，也注定要获得永生的复活，而不是永死，因为他们将在基督内复活？他们在基督内复活，因为他们属于基督的身体？他们属于基督的身体，因为基督是他们的头？\BibleRef{格前11:3}而基督是他们的头，因为只有一位中保在天主与人之间，就是人基督耶稣。除非通过祂的恩宠，他们信了祂的复活，否则祂不可能成为他们的中保。如果他们不知道祂将要降生成人，并且没有凭这信德正直而虔诚地生活，他们怎能相信祂的复活呢？如果基督的降生成人与他们无关，因为此事尚未发生，那么基督的审判也必然与我们无关，因为此事尚未发生。但是，如果我们凭着信德相信基督的审判（这审判尚未发生，但将要发生）而站在基督的右边，那么那些古圣祖就是凭信德相信祂的复活（这复活在他们时代尚未发生，但终将发生）而成为基督的肢体。\par

% structure: p0065 source=h2 target=subsection reason=relative-heading-rank; base=h2

\subsection{第三十二章 [XXVII.]——他用\Person{亚巴郎}的例子证明古圣祖曾相信基督的降生成人。}

因为不可认为那些古圣祖只受益于基督的天主性（它永远存在），而没有受益于祂的人性的启示（它尚未发生）。主耶稣说：\BibleQuote{亚巴郎渴望看见我的日子，他看见了，且高兴起来。}\BibleRef{若8:56}祂用短语\Emph{他的日子}来意指\Emph{他的时期}，这当然是明确的证据，证明\Person{亚巴郎}充满对基督降生成人的信仰。正是就此而言，祂有\Quote{时期}；因为祂的天主性超越一切时间，因为一切时间都是藉它创造的。然而，如果有人认为所讨论的短语必须理解为那个永恒的日子，它没有明天，也没有昨天——总之，就是祂与圣父同永恒的那永恒本身——那么\Person{亚巴郎}怎么会真正渴望这个呢，除非他知道将来属于祂的必死性，而祂的永恒是他所渴望的？或者，也许有人会将短语的含义局限于此，认为主所说的\Quote{他渴望看见我的日子}无非是\Quote{他渴望看见我}，我就是那无始无终的日子，或不灭的光，正如我们提到圣子的生命，关于这生命福音中说：\BibleQuote{圣父赐给圣子在自己内有生命。}\BibleRef{若5:26}这生命就是祂自己。所以，我们理解圣子自己就是生命，当祂说：\BibleQuote{我是道路、真理、生命}\BibleRef{若14:6}，关于祂也说过：\BibleQuote{祂是真天主，是永生。}\BibleRef{若一5:20}那么，假设\Person{亚巴郎}渴望看见圣子与圣父同等的天主性，而预先不知道祂将要在肉身中来临——如同某些哲学家寻求祂，他们对祂的肉身一无所知——那么\Person{亚巴郎}的另一个行为，当他吩咐仆人把手放在他的大腿下，并指着天上的天主起誓，\BibleRef{创24:2}难道还能被任何人理解为别样意思吗？这行为表明\Person{亚巴郎}深知天主将要在肉身中来临，而这肉身正是从他自己的大腿所出的后裔。\par

% structure: p0067 source=h2 target=subsection reason=relative-heading-rank; base=h2

\subsection{第三十三章 [XVIII.]——基督如何是我们的中保。}

关于这血肉与血，\Person{默基瑟德}在祝福\Person{亚巴郎}本人时，\BibleRef{创14:18}也给出了基督徒信徒非常熟悉的见证，以致很久以后在圣咏中向基督说：\BibleQuote{你照\Person{默基瑟德}的品位，永为司祭。}那时这还不是已成就的事实，而是未来的事；然而，祖先们的那信德（与我们的信德相同）曾咏唱它。现在，对于所有在亚当内找到死亡的人，基督有这样的效用：祂是生命的中保。然而，祂之所以是中保，不是因为祂与圣父平等；因为在这方面，祂自己与我们的距离如同圣父一样遥远；既然距离相同，又怎能有任何中介呢？因此，宗徒不是说他\Quote{只有一个中保在天主与人之间，就是耶稣基督}，而是说：\BibleQuote{人基督耶稣。}\BibleRef{弟前2:5}祂是中保，因为在祂是人——次于圣父，正如祂更接近我们，又高于我们，正如祂更接近圣父。这更明白地表达为：\BibleQuote{祂次于圣父，因为取了奴仆的形体；}\BibleRef{斐2:7}高于我们，因为无罪的玷污。\par

% structure: p0069 source=h2 target=subsection reason=relative-heading-rank; base=h2

\subsection{第三十四章 [XXIX.]——除基督外，无人曾得救。}

如今，凡是认为人性在任何时期都不需要第二亚当作其医师，因为它在第一亚当中并未败坏的人，都被证实为天主恩宠的敌人；并非在一个可与正统信仰相容的疑问或错误中，而是在那使我们成为基督徒的信仰准则本身中。那么，这些人在赞颂最初存在的人性时，为何说它远未受到恶习的污染？他们怎会忽视这一事实：当时的人们深陷于如此多不可容忍的罪恶中，以至于除了一位属神的人及其妻子，以及三个儿子和他们的妻子外，整个世界都在天主公义的审判中被洪水毁灭，如同后来索多玛的小地被火毁灭一样？因此，从那时起，当\BibleQuote{正如罪恶借着一人进入了世界，死亡借着罪恶也进入了世界；这样死亡就殃及了众人，因为众人都犯了罪}\BibleRef{罗 5:12}，我们整个本性的整体无疑已被败坏，落入了毁灭者的手中。除了通过救赎者的恩宠，从他手中没有人——确实没有一人——曾得解救，或正在被解救，或将要得解救。\par

% structure: p0071 source=h2 target=subsection reason=relative-heading-rank; base=h2

\subsection{第35章 [30]——为何在如此严厉的刑罚下命令婴孩受割损礼。}

圣经并未告诉我们，在亚巴郎之前，义人及其子女是否由任何身体或可见的标志所标记。亚巴郎本人确实接受了割损礼的标记，作为因信德而获得的正义的印证。\BibleRef{罗 4:11}他接受此礼时附有命令：从那时起，他家中所有的男婴在出生后第八天，在刚从母胎出来时，都要受割损礼；\BibleRef{创 17:10}这样，甚至那些尚不能用心相信以致成义的人，也应接受因信德而获得的正义的印证。这道命令附有如此可怕的惩罚，以致天主说：\BibleQuote{那在第八天包皮未受割损的男性，应从民间铲除。}\BibleRef{创 17:14}倘若追问如此重罚的正义性，那么这些人关于自由意志以及本性可嘉的健康与纯洁的全部论证，无论多么巧妙，岂不会被击碎、打垮并粉碎成原子？因为，请告诉我，一个婴孩凭自己的意志做了什么恶事，以至于因他人未给他行割损礼的疏忽，他自己必须被定罪，且受如此严厉的定罪，以致他的灵魂应从其人民中被铲除？这可怕之处并非针对任何暂时的死亡，因为义人死去时，通常反而是说：\BibleQuote{他便归到他亲族那里去了}\BibleRef{创 25:17}，或\BibleQuote{他便归到他祖宗那里去了}\BibleRef{加上 2:69}：因为当一个人自己的民众本身就是天主的子民时，将他与他的民众分离并不长久令人恐惧。\par

% structure: p0073 source=h2 target=subsection reason=relative-heading-rank; base=h2

\subsection{第36章 [31]——驳斥柏拉图主义者关于灵魂在身体前已存在的意见。}

那么，如此严厉的定罪有何意义，既然没有犯任何有意的罪？因为并不像某些柏拉图主义者所认为的那样，每个这样的婴孩在其灵魂中因它在今世生活之前凭自己意志所做的事而受此报应，仿佛它在现有身体状态之前就拥有自由选择过好或过坏的生活；因为宗徒保禄说得很清楚，他们在出生之前既没有行善也没有行恶。\BibleRef{罗 9:11}因此，一个婴孩凭什么理由应如此毁灭性地受惩罚，除非因为他属于该毁灭的群体，并正当地被视为因亚当而生，在古债的束缚下被定罪，除非他非因债而是因恩宠从束缚中被释放？而除了通过我们的主耶稣基督，这恩宠还能是什么？古代犹太人不仅在其他神圣制度中，而且在他们对包皮的割损礼中，无疑已包含了对祂来临的预兆。因为第八天，在星期的循环中，成为主的日子，主在这一天从死者中复活；基督是磐石\BibleRef{格前 10:4}，从中形成了割损的石刀；\BibleRef{出 4:25}包皮之肉便是罪恶的身体。\par

% structure: p0075 source=h2 target=subsection reason=relative-heading-rank; base=h2

\subsection{第37章 [32]——基督在何种意义上被称为\Quote{罪}。}

圣事礼仪在祂的来临之后有所改变，因为祂的来临是这些礼仪所预表的；但中保的帮助并无改变，祂甚至在其降生成人之前，一直藉着古时祂身体的肢体对祂降生成人的信德来拯救他们；至于我们自身，虽然我们死于罪恶和未受割损的肉身上，却与基督一同活过来了，在祂内我们受了割损，不是人手所行的割损，而是由旧日人手割损所预表的割损，好使那从亚当遗传给我们的罪恶的身体被解除\BibleRef{罗 6:6}。被定罪的来源的繁衍定我们的罪，除非我们被相似有罪的肉身的形像所洁净，而祂就是在这样的肉身中被派遣，自己无罪，却为我们的缘故成了罪，在肉身上定了罪的罪。因此宗徒说：\BibleQuote{我们在基督的位份上恳求你们，与天主和好吧！因为祂使那无罪的，替我们成了罪，好叫我们在祂内成为天主的正义。}\BibleRef{格后 5:20}因此，我们与祂和好的天主，使祂为我们成了罪——就是说，成为祭献，藉此我们的罪得以赦免；因为\Quote{罪}字也指赎罪祭。的确，祂是为我们的罪被祭献的，在人类中只有祂是无罪的，正如在古时，人们从羊群中寻找一只来预表那将要来临医治我们过犯的无缺者。因此，无论婴孩在出生后哪一天受洗，他都如同在第八日受了割损；因为他在祂内受了割损，祂被钉十字架后第三天复活，但按周算是第八天。他受割损是为了脱去罪恶的身体；换言之，属神重生的恩宠要消除因肉身生育的传染所欠下的债务。因为\BibleQuote{没有一人是洁净而不染污秽的}（请问是什么污秽，不是罪的污秽吗？），\BibleQuote{连在世只活了一天的婴孩也不洁净。}\BibleRef{约 14:4}\par

% structure: p0077 source=h2 target=subsection reason=relative-heading-rank; base=h2

\subsection{第三十八章——原罪并不使婚姻成为恶事}

但他们这样争论说：\Quote{难道婚姻不是恶事，由婚姻所生的人不是天主的化工吗？}好像婚姻生活的善在于那种私欲偏情，那些不认识天主的人以此爱他们的妻子——这是宗徒所禁止的\BibleRef{得前 4:5}；而不在于那夫妇间的贞洁，藉此肉欲被导向为生育子女而命定的善的目的。或者好像，一个人若不是天主的化工，不仅生于合法婚姻，甚至生于奸淫或通奸也不可能。但在当前的问题中，当问题不在于为了什么需要一位造物主，而是为了什么需要一位救主时，我们不必考虑自然生育有什么善，而是要考虑罪恶有什么恶，我们的本性确实因罪恶而败坏。无疑，两者同时生成——本性与本性的败坏；一个是善，另一个是恶。一个来自造物主的恩赐，另一个来自我们起源的定罪；一个的原因在于至上天主的善意，另一个的原因在于第一人的败坏意志；一个显示天主为受造物的\Emph{造主}，另一个显示天主为不顺服的惩罚者：简言之，同一个基督，为创造那一个成了人的\Emph{造主}，并为医治那另一个\Emph{成了人}。\par

% structure: p0079 source=h2 target=subsection reason=relative-heading-rank; base=h2

\subsection{第三十九章——婚姻中三件善事与可赞美之事}

因此，婚姻在一切属于婚姻状态的事物中是一种善。这些事物有三：它是生育的有序方式，它是贞洁的保证，它是结合的纽带。就其生育的序次而言，经上说：\BibleQuote{所以我愿意年轻的寡妇出嫁，生育儿女，管理家务}\BibleRef{弟前 5:14}；就其贞洁的保证而言，论及婚姻说：\BibleQuote{妻子对自己的身体没有主权，而是丈夫有；同样，丈夫对自己的身体也没有主权，而是妻子有}\BibleRef{格前 7:4}；就其作为结合的纽带而言：\BibleQuote{天主所结合的，人不可拆散}\BibleRef{玛 19:6}。关于这些要点，我们并未忘记，在主赐给我们的才能中，我们已在你们所熟悉的其他著作中作了充分的论述。关于这一切，经上有一段普遍的赞美：\BibleQuote{婚姻在各方面都应受尊重，床第也不可玷污}\BibleRef{希 13:4}。因为婚姻状态既为善，它便在私欲偏情的恶中产生极大的益处；因为不是情欲，而是理性，善用了私欲偏情。情欲在于那\Quote{违命的}肢体的法律，宗徒称之为\BibleQuote{反抗心智的法律}\BibleRef{罗 7:23}；而理性则在于婚姻状态的法律，它善用了私欲偏情。然而，如果任何善不可能从恶中产生，那么天主就不能从奸淫的拥抱中创造人。因此，每当人在奸淫中诞生时，那该受罚的奸淫之恶并不归咎于天主，因为天主在人的恶行中确实施行了善工；同样，那使人在肢体叛乱中感到羞愧的一切——这叛乱曾带来羞红的谴责，使人在犯罪后用无花果树叶遮盖这些肢体\BibleRef{创 3:7}——并不归咎于婚姻，凭借婚姻，夫妻的拥抱不仅是允许的，甚至是有益和尊贵的；而是归咎于那悖逆之罪，罪罚是人发现自己的肢体仿效他自己对天主所犯的悖逆，与自己作对。于是，人对自己的行为感到羞惭，因为这些肢体不再随理性意志的指挥而行动，而是仿佛随心所欲，受情欲的煽动；他便想出了遮盖物，以隐藏那些他认为值得羞耻的肢体。因为人作为天主的化工，本不应蒙羞；那造物主认为适合形成和指定的肢体，绝无意使受造物脸红。因此，那简单的裸体既不令天主不悦，也不令人羞惭：没有什么可羞耻的，因为起初并未发生任何应受惩罚的事。\par

% structure: p0081 source=h2 target=subsection reason=relative-heading-rank; base=h2

\subsection{第四十章【第三十五题】——婚姻在犯罪之前已经存在。天主的祝福如何在我们的原祖父母身上运作。}

然而，即使罪尚未存在，婚姻无疑已经存在；并且为此原因，女人被造为男人的助手，而不是另一个男人。此外，天主的话：\BibleQuote{你们要生育繁殖}\BibleRef{创 1:28}并非预言要受谴责的罪，而是对婚姻生育力的祝福。因为借着这些祂不可言喻的话——我是指那蕴含于祂智慧真理中的神圣方法，万物藉以造成——天主赋予原祖夫妇生育的能力。然而，假如本性未曾因罪而受辱，我们万不可认为乐园中的婚姻会是这样：在其中，生育的肢体是受情欲的炽热所激动，而非受意志的命令以产生后代——如同脚为行走，手为劳动，舌为言语一样。而且，不会像现在这样，童贞的贞洁因着浑浊的热力而被破坏以怀孕，而是会服从于最温柔的爱的能力；如此，交合中的女子没有疼痛，没有流血，分娩的母亲也没有呻吟。然而，人因我们的有死状态的实际经验中没有证实这一点，便拒绝相信。本性因罪而受损，从未体验过那原始纯洁的例子。但我们是对有信德的人说话，他们已学会相信默感的圣经，即使没有实际现实的例子。因为我现在怎能用实际来\Emph{证明}一个人是由尘土所造，没有父母，并且一个妻子是由他的肋骨为他形成的呢？\BibleRef{创 2:7}然而，信德相信眼睛不再发现的事。\par

% structure: p0083 source=h2 target=subsection reason=relative-heading-rank; base=h2

\subsection{第四十一章【第三十六题】——私欲和产痛来自罪。我们的肢体何时成为羞耻的原因。}

因此，我们虽然无法证明，那原始婚姻中的婚姻行为是在没有情欲骚动的情况下平静完成的，并且生殖器官的运动，如同身体其他肢体一样，并非受情欲热忱的驱使，而是受意志选择的引导（倘若没有罪的耻辱介入，这种情况在婚姻中本会持续），然而，根据圣经中所记载的神圣权威，我们有充分理由相信，这就是婚姻生活的原始状态。尽管确实，我并未被告知婚姻拥抱没有伴随着色欲；正如我也未发现记载说分娩没有伴随着呻吟与痛苦，或者实际出生之后不会导致死亡；但同时，如果我遵循圣经的真理，母亲的产痛和人类后裔的死亡，若非罪已先行，是绝不会发生的。那使男人和女人遮盖腰际的事也不会发生；因为在同一神圣记录中明确写道，罪先被犯下，然后紧接着就是这种羞耻的隐藏。\BibleRef{创 3:7} 因为，除非某种不雅的动态向他们的眼睛——当然不是闭着的，但并未注意到这一点，即未留意——宣告那些肢体需要被纠正，否则他们不会在自己身上察觉到任何值得羞耻或需要隐藏的事，而天主所造的一切本是全然值得赞美的。的确，如果罪没有先行发生，他们不会在悖逆中胆敢犯罪，也就不会随后出现羞耻所要隐藏的耻辱。\par

% structure: p0085 source=h2 target=subsection reason=relative-heading-rank; base=h2

\subsection{第42章 [XXXVII.]— 情欲之恶不应归咎于婚姻。婚姻法令的三项善果：子孙、贞洁与圣事性结合。}

显然，那不可归咎于婚姻，即使没有它，婚姻仍然存在。婚姻之善并不因恶而被除去，尽管恶被婚姻用于善的目的。然而，这是现在有死之人的状况，使婚姻交合与情欲同时发生；因此，那些不愿或不能区分二者的人，便认为不仅情欲该受谴责，而且连婚姻行为本身，尽管合法且尊荣，也当受责备。他们也不留意于婚姻状态之善，即婚姻的荣耀：我指的是子孙、贞洁与盟约。但连婚姻也为之羞耻的恶，并非婚姻的过失，而是肉体的情欲。然而，因为没有这恶就无法实现婚姻的善的目的，即生育子女，每当接近这一过程时，人们寻求隐秘，去除见证人，甚至避免那些由此过程出生的子女在达到能观察的年龄时在场。于是，婚姻被允许实现其状态中一切合法之事，只是不可忘记隐藏一切不当之事。由此得出，婴儿虽不能犯罪，却仍在罪的传染中出生——并非因合法之事，而是因不当之事：因为从合法之事出生的是本性，从不当之事出生的是罪。如此出生的本性，其创造者是造人的天主，祂在婚姻法律下结合男女；但罪的作者是欺骗人的魔鬼的诡诈和同意之人的意志。\par

% structure: p0087 source=h2 target=subsection reason=relative-heading-rank; base=h2

\subsection{第43章 [XXXVIII.]— 人类的后裔，甚至在出生之前，已在根源上被判罪。仅为快感而进行的婚姻行为并非没有小罪。}

在那里，天主除了以公义的判决定那故意犯罪的人连同其苗裔有罪之外，别无作为；当然，所有尚未出生的人也公义地在有罪的根源中被定罪。在这被定罪的根源中，肉身生育掌握着每个人；唯有属灵的再生才能释放他。因此，对于重生的父母，如果他们继续处于同样的恩宠状态，由于他们已经获得的罪的赦免，这无疑不会产生有害的后果，除非他们滥用恩宠——不仅包括各种不法淫乱，甚至在婚姻状态中，每当夫妻进行生育时，并非出于自然繁衍种族的愿望，而是纯粹因放纵而沉溺于情欲的满足。宗徒允许夫妻\BibleQuote{不要彼此亏负，除非两相情愿，暂时分房，为专务祈祷}\BibleRef{格前 7:5}，他是以宽容的许可而非命令来让步；但这种让步的形式本身显然暗示着一些过失。然而，婚姻契约所指的旨在生育子女的婚姻拥抱，就其本身并单纯考虑，不涉及邪淫，是善而正当的；因为，虽然由于这死亡的身体（尚未因复活而更新）的缘故，没有一定程度的兽性运动就无法进行，这种运动使人类感到羞耻，但这拥抱本身并非罪，当理性将私欲偏情用于善的目的，而不被其征服向恶时。\par

% structure: p0089 source=h2 target=subsection reason=relative-heading-rank; base=h2

\subsection{第44章 [XXXIX.]— 就连重生者的子女也是在罪中出生。圣洗圣事的效果。}

肉身的这种贪欲，只要它存在于我们内，原本就会造成损害——若非罪的赦免如此有益，以致它在人身上（无论是生来的人，还是重生的人）都存在，但在前者中是既有害又存在，而在后者中仅仅存在却从无害处。在未重生的人中，它的确有害到如此程度：除非他们重生，否则即使出生于重生的父母，对他们也毫无益处。我们本性的缺陷如此深刻地印在我们的后代身上，使其成为有罪的，即使这同一缺陷的罪过在父母身上已藉着罪的赦免而被洗去——直到在最后的重生中，一切因人类意志的同意而终结于罪过的缺陷都被耗尽并消除。这最后的重生将与肉身的未来复活完全相同；那时我们不仅不犯罪，甚至摆脱那些因我们向其屈从而引我们犯罪的败坏欲望。我们甚至现在就已通过那圣洗之水的恩宠向着这有福的圆满迈进。那如今更新我们灵魂的同一重生，使我们过去的一切罪过都得以赦免；日后，如所期待的那样，也将作用于那同一肉身，使其更新而进入永生；藉着复活到一个不朽的状态，一切罪过的诱因都将从我们的本性中被清除。但这救恩如今仅在望德中完成，尚未在事实上实现；它不是现世拥有的，而是以忍耐期待的。\newline{}【XL】因此，在同一个圣洗的水泉中，有一种完整而彻底的洁净，不仅包括现今在洗礼中被赦免的一切罪过——这些罪过因我们同意邪恶欲望及其行为而成罪——也包括这些邪恶欲望本身；这些欲望若未经我们同意，便不构成罪过；它们虽在现世生活中未被移除，但在来世将不复存在。\par

% structure: p0091 source=h2 target=subsection reason=relative-heading-rank; base=h2

\subsection{第四十五章——人的得救与其被俘的特征相称。}

因此，我们所谈的这种败坏之罪，将留在重生者的肉身后裔中，直到在他们身上也藉着重生的洗礼之水被洗去。一个重生的人不重生后代，而是按肉身生育儿子；这样，他传给后代的不是重生的状态，而仅仅是出生的状态。因此，无论一个人是不信者，还是完全的信仰者，他都不会生下有信德的儿女，而是生下有罪的儿女；同样，无论是野橄榄的种子，还是栽培橄榄的种子，产生的都不是栽培橄榄，而是野橄榄。同样，一个人的第一次出生使他处于那种束缚中，只有第二次出生才能解救他。魔鬼掌控他，基督释放他；欺骗厄娃的那位掌控他，玛利亚的儿子释放他；通过女人接近男人的那位掌控他，从未接近男人而由女人所生的那位释放他；向女人注入情欲之因的那位掌控他，毫无情欲地在女人内受孕的那位释放他。前者能通过一个人掌控所有人；除了他不曾掌控的那一位，没有人能将他们从他的权力中解救出来。事实上，教会的圣事本身——教会依据极古老的传统权威，以应有的仪式施行这些圣事（这些人尽管认为圣事在婴儿身上是模仿性的而非真实性的，但仍不敢公开拒绝它们）——我说，圣教会的事本身就足以表明，即使刚出母腹的婴儿，也是藉着基督的恩宠从魔鬼的束缚中得救的。因为，且不说他们为了赦罪而受洗这一并非虚假而是真实可信的奥迹，在这之前，还有对敌对力量的驱魔和呵斥；他们藉着那些带领他们前来受洗之人的口，宣认弃绝这力量。藉着这些对隐藏之现实的圣化而明显的标志，他们被显示为从最坏的压迫者转向最卓越的救赎者；祂为我们承担了我们的软弱，捆绑了那壮士，为要夺取他的器物；\BibleRef{玛 12:29} 因为天主的软弱比人甚至比天使都更强。因此，当天主解救小人物和大人物时，祂在两种情况下都表明，宗徒是受真理的指引而说话的。因为祂不仅解救成年人，也解救小婴孩，将他们从黑暗的权势中救出，以便把他们移入祂爱子的国度。\BibleRef{哥 1:13}\par

% structure: p0093 source=h2 target=subsection reason=relative-heading-rank; base=h2

\subsection{第四十六章——相信原罪的困难。人的恶是野兽的本性。}

谁都不应感到惊讶并问：\Quote{为什么天主的美善要为魔鬼的恶毒创造任何东西来占据？}事实是，天主的恩赐是以同样的慷慨赐予祂受造物的生殖元素，祂用这慷慨\BibleQuote{叫日头照恶人，也照善人；降雨给义人，也给不义的人。}\BibleRef{玛 5:45}祂就是以如此大的慷慨祝福了种子，并藉着祝福而构成了它们。这祝福并没有因使我们陷于定罪的过失而从我们卓越的本性中被消除。的确，由于惩罚人的天主的公义，这致命的缺陷已经如此得胜，以致人生来就带有原罪的过失；但它的影响尚未扩展到阻止人的出生。在成年人的身上也是如此：无论他们犯什么罪，都不能从人身上消除他的人性；相反，天主的工程仍然是好的，尽管恶人的行为是邪恶的。因为虽然\Quote{人躺在尊位上而不居住；由于无知，与野兽相似，成了跟它们一样的，}但相似并非绝对到使他成为野兽。两者之间无疑有比较，但不是由于本性，而是由于恶——不是在野兽中的恶，而是在本性中的恶。因为人与野兽相比是如此卓越，以致人的恶成了野兽的本性；然而人的本性从未因此改变成野兽的本性。所以，天主定罪人是因为使他的本性蒙受耻辱的过失，而不是因为他的本性，本性并没有因它的过失而被毁灭。但愿我们不要认为野兽应受定罪的判决！它们理应免于我们的痛苦，因为它们不能分享我们的幸福。那么，人服从于不洁之神——不是由于本性，而是由于他在出生污点中所沾染的、并非来自天主的工程而是来自人的意志的不洁——这有什么奇怪或不公呢？因为那不洁之神本身作为神是好的，但作为不洁则是邪恶的。因为前者来自天主，是祂的工程，而后者来自人的意志。因此，更强的本性，即天使的本性，由于恶与后者的结合，使较低的即人的本性处于从属地位。因此，比天使更强的中保，为了人的缘故而成了软弱的。\BibleRef{格后 8:9}如此，毁灭者的骄傲被救主的谦卑所摧毁；那以天使的力量向人类夸耀的，被圣子在祂所取的人性软弱中击败了。\par

% structure: p0095 source=h2 target=subsection reason=relative-heading-rank; base=h2

\subsection{第四十七章——\Person{安波罗削}支持原罪的语句}

现在，我们即将结束本书，以为在此关于\Emph{原罪}的题目上，我们应当像先前在《\Emph{论恩宠}》的论著中所做的一样——引证天主的那位仆人、总主教\Person{安波罗削}作为反对这些人的有害言论的证据，\Person{白拉奇}声称这位圣者的信德在拉丁教会的作者中是最完美的；因为\Emph{恩宠}在消除\Emph{原罪}上尤其被尊崇。在圣\Person{安波罗削}所写的《\WorkTitle{论复活}》中，他说：\Quote{我堕落在亚当内，在亚当内被逐出乐园，在亚当内死亡；除非祂在亚当内找到我，祂不会召回我——以致正如我在他身上有罪债，且该死，我也在基督内成义。}然后，在驳斥诺瓦替亚努派时，他又说：\Quote{我们众人都是在罪中出生的；我们的起源就在于罪；正如你读到\Person{达味}所说：\BibleQuote{我是在罪孽里生的，我母亲在罪中怀了我。}因此保禄的肉体是\BibleQuote{死亡之身}；\BibleRef{罗 7:24}正如他亲自说：\BibleQuote{谁能救我脱离这死亡之身呢？}基督的肉却定罪了罪，这罪祂并未因出生而经历，且藉死亡而钉死了它，好使我们的肉体藉恩宠获得成义，在此之前则因罪而不洁。}同一圣人在他的《\WorkTitle{依撒意亚注释}》中，论及基督时说：\Quote{因此，祂作为人，在一切事上受过试探，在人的相似内忍受了一切；但因圣神而生，祂无罪可犯。因为每个人都是说谎者，除了天主以外，没有无罪者。因此，这是一个观察而确定的事实：凡由男女结合而生的人，都不能免于罪。的确，凡免于罪的，也免于这样的孕育和出生。}此外，在解释路加福音时，他说：\Quote{并非与丈夫的同房开启了童贞女之胎的秘密；而是圣神将无玷的种子注入她未染的胎中。因为主耶稣独自一人由妇女所生是圣洁的，因祂未曾接触世俗的腐败，由于祂无玷出生的新颖；祂反而以天上的尊威将其驱逐。}\par

% structure: p0097 source=h2 target=subsection reason=relative-heading-rank; base=h2

\subsection{第四十八章——\Person{白拉奇}被\Person{安波罗削}正确谴责并真正反对}

然而，这位属于天主的人所说的话，却遭到了\Person{白拉奇}的反驳，尽管他对该作者不吝赞美之词，因为他自己宣称：\Quote{我们被生出时，既无德行，亦无恶习}。那么，剩下还有什么呢？除非\Person{白拉奇}谴责并弃绝他自己的这个错误，否则他就得对自己这样引用\Person{盎博罗削}表示遗憾。然而，既然真福\Person{盎博罗削}，作为一位公教主教，在上述引文中是按照公教信德来表达的，那么，\Person{白拉奇}连同他的门徒\Person{塞莱斯提乌}，因背离了信德的真道而被天主教会权威公正地定罪，因为他未曾悔改曾对\Person{盎博罗削}大加赞赏却同时持有与他对立的观点。我十分清楚，你们怀着怎样的如饥似渴之情阅读任何为建立并巩固信德而写成的文字；但是，尽管它对于达到这一目的有助益，我还是必须最后结束这篇论文。\par

\SourceRecord{Translated by Peter Holmes and Robert Ernest Wallis, and revised by Benjamin B. Warfield.；Nicene and Post-Nicene Fathers, First Series；Vol. 5.；Edited by Philip Schaff.；Buffalo, NY: Christian Literature Publishing Co.,；1887.；<http://www.newadvent.org/fathers/15062.htm>.}
